\chapter[Gauge-Retained Network State and Fixed Sectors]{Gauge-Retained Network State\\and Fixed Sectors}
\label{ch:network-state}

Grouping the agents at $c_*$ into blocks does not average away their gauge data: links are independent
group elements and closed walks carry holonomy. This chapter asks what a coarse description of a block
must retain of that data and which part of the block's state survives exact aggregation. A rooted spanning
tree retains tree transports, the based holonomy, and block-dressed boundary edges as one root-framed datum
(\Cref{prop:rg-rooted-datum}). The fixed sector of a block is the kernel of its connection Laplacian, with
positive-definite weights the holonomy-fixed root values (\Cref{prop:cg-kernel-holonomy}), and transported
marginal laws admit a parent law exactly when they occupy one holonomy-stabilized mode
(\Cref{thm:cg-holonomy-kl-marginal}). Cut links compress exactly into rectangular endpoint maps
(\Cref{thm:cg-rectangular-endpoint-closure}), and into one group-valued link with new self terms exactly
when an explicit matrix is invertible (\Cref{lem:cg-cut-completed-square}), with one balanced choice of
link and weight (\Cref{prop:cg-balanced-cut-presentation}) that no presentation at the level of laws
reproduces once the cut has excess (\Cref{prop:cg-no-balanced-law-presentation}) and that admits no continuous
extension to all singular cuts (\Cref{open:cg-cut-compression}); the coarse frame is invisible to
the constituents and cannot be computed from their frames by a symmetric equivariant rule
(\Cref{prop:cg-frame-sector}); and a record at $c_*$ distinguishes loop holonomies at the level of laws
(\Cref{thm:cg-loop-operational-holonomy}). These results supply the retained gauge data of
\Cref{ch:agent-network-rg,ch:full-graph-meta-agent-vfe}, the formation witnesses of
\Cref{ch:finite-depth-formation}, and the testability discussion of \Cref{ch:philosophy}; Supplement G
realizes them for Gaussian families.

\paragraph{Scope.} The fixed-sector rank is an admissibility criterion: it assigns no partition cost,
keeps singletons admissible, and does not require constituent beliefs to coincide. The marginal-law score
takes its block and weights as inputs, selects no partition, supplies no Markov kernel, and does not
recover the dependence structure of a joint law. A cross map parallel for the links neither identifies the
belief and model bundles nor orders their two scores, and the operational record, which sees a loop only
through its represented holonomy at the probe, recovers no link values. \status{NOT-CLAIMED} The obstruction
to a continuous extension of balanced cut links is \Cref{open:cg-cut-compression}.

\section{Levelwise fiber state and scale intertwiners at \texorpdfstring{$c_*$}{c*}}
\label{sec:rg-common-principal-state}

\definitionheading{Levelwise fiber state}{def:rg-geometric-state}
At depth $\ell$ fix a Lie group $G_\ell$ and a right $G_\ell$-torsor $\mathscr P_\ell$, the fiber over $c_*$
of the principal bundle of that depth, containing one frame of each agent or meta-agent of depth $\ell$,
used for both channels. Law fibers $\mathcal B_{\ell,x}$, $x\in\{b,m\}$, carry left actions
$\widehat\rho_{\ell,x}$ of $G_\ell$, which may be inequivalent, with associated fibers
$\mathcal E_{\ell,x}=\mathscr P_\ell\times_{\widehat\rho_{\ell,x}}\mathcal B_{\ell,x}$. The depth-$\ell$
interaction multigraph $\mathfrak I_\ell$ carries links $\Theta^x_e\in G_\ell$ with
$\Theta^x_{\bar e}=(\Theta^x_e)^{-1}$ (\Cref{def:geo-graph-links}), every block $I$ of depth $\ell$ carries
its based holonomy groups $\mathrm{Hol}^x_I$ (\Cref{def:rg-rooted-tree-data}), and declared cross maps
$\Phi_\ell:\mathcal E_{\ell,b}\to\mathcal E_{\ell,m}$ and $\widetilde\Phi_\ell:\mathcal E_{\ell,m}\to\mathcal E_{\ell,b}$
need not be invertible or mutually inverse. The fiber state of depth $\ell$ is
\begin{equation}
\mathfrak X_\ell=\Bigl(G_\ell,\mathscr P_\ell;\ \mathcal E_{\ell,b},\mathcal E_{\ell,m};\
\mathfrak I_\ell,(\Theta^b_e,\Theta^m_e)_e,(\mathrm{Hol}^b_I,\mathrm{Hol}^m_I)_I;\
\Phi_\ell,\widetilde\Phi_\ell;\ \ldots\Bigr),
\label{eq:rg-full-state}
\end{equation}
where the ellipsis stands for law, kernel, or operator components declared by a realization.

\propositionheading{Scale maps and intertwiners at \texorpdfstring{$c_*$}{c*}}{prop:rg-scale-intertwiner}
(i) A map $\Phi_\ell:\mathcal E_{\ell,b}\to\mathcal E_{\ell,m}$, expressed in a frame $u$ by the fiber map
$f_u$ with $\Phi_\ell[u,w]=[u,f_u(w)]$, has one expression $f^{bm}_\ell$ in every frame exactly when
$f^{bm}_\ell\circ\widehat\rho_{\ell,b}(g)=\widehat\rho_{\ell,m}(g)\circ f^{bm}_\ell$ for every $g\in G_\ell$,
which is also the condition that $\Phi_\ell$ commutes with every automorphism of $\mathscr P_\ell$ and
constrains the covariance of a cross map, not its existence; the reverse direction is analogous. For a
homomorphism $\kappa_\ell:G_\ell\to G_{\ell+1}$, points
$p_0\in\mathscr P_\ell$ and $p_1\in\mathscr P_{\ell+1}$, and fiber maps
$q_{\ell,x}:\mathcal B_{\ell,x}\to\mathcal B_{\ell+1,x}$, $x\in\{b,m\}$, consider
\begin{align}
\mathsf P_\ell(p\cdot g)&=\mathsf P_\ell(p)\cdot\kappa_\ell(g),
\label{eq:rg-principal-scale-map}\\
\widehat q_{\ell,x}[p,w]&=[\mathsf P_\ell(p),q_{\ell,x}(w)],
\label{eq:rg-associated-scale-map}\\
q_{\ell,x}\circ\widehat\rho_{\ell,x}(g)&=\widehat\rho_{\ell+1,x}(\kappa_\ell(g))\circ q_{\ell,x}
\qquad\text{for every }g\in G_\ell.
\label{eq:rg-scale-intertwiner}
\end{align}
(ii) Exactly one map $\mathsf P_\ell:\mathscr P_\ell\to\mathscr P_{\ell+1}$ satisfies $\mathsf P_\ell(p_0)=p_1$
and \eqref{eq:rg-principal-scale-map}. (iii) The formula \eqref{eq:rg-associated-scale-map} defines a map
$\widehat q_{\ell,x}:\mathcal E_{\ell,x}\to\mathcal E_{\ell+1,x}$ exactly when \eqref{eq:rg-scale-intertwiner}
holds. (iv) For composable steps the homomorphisms, torsor maps, fiber maps, and induced maps compose, the
rightmost factor acting first, and on law fibers given by pushforward the composite is the pushforward of
the composed kernel, by \eqref{eq:cg-law-functor}. \status{ESTABLISHED}

\noindent\emph{Proof.}
(i) Each point of $\mathcal E_{\ell,b}$ has one representative with first entry $u$, and
$[u\cdot g,w]=[u,\widehat\rho_{\ell,b}(g)w]$ gives
$f_{u\cdot g}=\widehat\rho_{\ell,m}(g)^{-1}\circ f_u\circ\widehat\rho_{\ell,b}(g)$, as in
\eqref{eq:geo-phi-gauge-law}. An automorphism of the torsor, $u\mapsto u\cdot\kappa$, acts by
$[u,w]\mapsto[u,\widehat\rho(\kappa)w]$, so commuting with all of them is the same condition. (ii) The
action is free and transitive, so $\mathsf P_\ell(p_0\cdot g):=p_1\cdot\kappa_\ell(g)$ is well defined,
satisfies \eqref{eq:rg-principal-scale-map}, and is forced by it. (iii) The representatives $(p,w)$ and
$(p\cdot g,\widehat\rho_{\ell,x}(g)^{-1}w)$ of one point have images $[\mathsf P_\ell(p),q_{\ell,x}(w)]$ and
$[\mathsf P_\ell(p),\widehat\rho_{\ell+1,x}(\kappa_\ell(g))q_{\ell,x}(\widehat\rho_{\ell,x}(g)^{-1}w)]$, which
agree for every $(p,w,g)$ exactly under \eqref{eq:rg-scale-intertwiner}. (iv) The composites satisfy the
defining identities by substitution. $\square$

On a chain of depths these composites form a scale diagram (\Cref{def:rg-scale-diagram,prop:rg-scale-functoriality});
at the information-geometric tier an induced map also needs the pushed laws to stay in the coarse model and
$q_{\ell,x}$ to be smooth. For an edge copy $e:j\to i$ of $\mathfrak I_\ell$, the cross-scale morphism
defects and the link-intertwining defects are the ordered pairs
\begin{equation}
\begin{aligned}
\mathbf D^\Phi_\ell&=\bigl(\widehat q_{\ell,m}\circ\Phi_\ell,\ \Phi_{\ell+1}\circ\widehat q_{\ell,b}\bigr),
&
\mathbf D^{\widetilde\Phi}_\ell&=\bigl(\widehat q_{\ell,b}\circ\widetilde\Phi_\ell,\ \widetilde\Phi_{\ell+1}\circ\widehat q_{\ell,m}\bigr),\\
\mathbf D^\Phi_e&=\bigl(\widehat{\mathsf t}^m_e\circ\Phi_\ell,\ \Phi_\ell\circ\widehat{\mathsf t}^b_e\bigr),
&
\mathbf D^{\widetilde\Phi}_e&=\bigl(\widehat{\mathsf t}^b_e\circ\widetilde\Phi_\ell,\ \widetilde\Phi_\ell\circ\widehat{\mathsf t}^m_e\bigr),
\end{aligned}
\label{eq:rg-cross-morphism-defects}
\end{equation}
where $\widehat{\mathsf t}^x_e[p,w]=[\mathsf t^x_e(p),w]$ is induced by the link map $\mathsf t^x_e$ of
\Cref{prop:geo-edge-twist-invariance}. A defect vanishes when its two components are equal; subtracting
them needs a vector or affine realization.

\propositionheading{Covariance of the defect diagrams}{prop:rg-defect-covariance}
The maps $\widehat{\mathsf t}^x_e$ are well defined and act from the frame of $j$ to the frame of $i$ as
$\widehat\rho_x(\Theta^x_e)$. Under passive re-expressions at depths $\ell$ and $\ell+1$, both components
of every pair in \eqref{eq:rg-cross-morphism-defects} acquire the same coordinate changes at source and
target, so the equality locus of each pair is gauge invariant; in a fiberwise-linear realization the
difference of the components transforms by the same left and right laws. \status{ESTABLISHED}

\noindent\emph{Proof.}
Right equivariance of $\mathsf t^x_e$ gives
$[\mathsf t^x_e(p\cdot g),\widehat\rho_x(g)^{-1}w]=[\mathsf t^x_e(p),w]$, and
$\mathsf t^x_e(u_j\cdot g)=u_i\cdot\Theta^x_eg$ gives the coordinate form. Both components of each pair are
frame-free maps between the same two associated fibers, so their frame expressions are conjugated by the
same coordinate changes, and equality of two maps does not depend on the frames. $\square$

\section{Rooted trees, based holonomy, and dressed edges}
\label{sec:rg-gauge-cross-scale}

Fix a depth and suppress its index. Let $\mathcal P$ be a partition of the agents into blocks and
$\mathfrak I_I$ the internal multigraph of a block $I$, with its loops and parallel copies. A passive
re-expression $g=(g_i)$ acts by $\Theta^x_e\mapsto g_i\Theta^x_eg_j^{-1}$ for $e:j\to i$
\eqref{eq:geo-regime-two-gauge-law} and by $z^x_i\mapsto\rho_x(g_i)z^x_i$ on sample coordinates
\eqref{eq:geo-represented-frame-change}.

\definitionheading{Rooted spanning-tree data}{def:rg-rooted-tree-data}
The component-indexed meta-agent set of $\mathcal P$ is
\begin{equation}
\widehat{\mathcal P}=\{(I,\mathsf c):I\in\mathcal P,\ \mathsf c\in\pi_0(\mathfrak I_I)\}.
\label{eq:rg-component-meta-index}
\end{equation}
From here on $I,J$ denote elements of $\widehat{\mathcal P}$, each connected with one root $r_I$; the state
of a disconnected block is the product over its components, and collapsing that product to one root, row,
or feature is an additional coarse channel. Choose a rooted spanning tree in each $I$, nested across
depths. With the path products $\Theta^x_p$ and holonomies $\mathrm{Hol}^x(p)$ of
\eqref{eq:geo-link-holonomy}, the tree transport of $i\in I$ along the tree path $p_{I,i}$ from $i$ to
$r_I$, the holonomy of the fundamental cycle of a non-tree copy $e:j\to i$ of $\mathfrak I_I$, the based
holonomy group of $I$ at a vertex $r$, and the block-dressed edge of a boundary copy $e:j\to i$ from $J$ to
$I$ are
\begin{align}
&\Theta^x_{I\leftarrow i}=\Theta^x_{p_{I,i}},
\qquad
\mathrm{Hol}^x(p_e)=\Theta^x_{I\leftarrow i}\Theta^x_e(\Theta^x_{I\leftarrow j})^{-1},
\label{eq:full-graph-tree-transport-law}\\
&\mathrm{Hol}^x_I(r)=\{\Theta^x_p:p\text{ a closed walk of }\mathfrak I_I\text{ at }r\}\subseteq G,
\qquad\mathrm{Hol}^x_I=\mathrm{Hol}^x_I(r_I),
\label{eq:full-graph-based-holonomy}\\
&\widehat\Theta^x_e=\Theta^x_{I\leftarrow i}\Theta^x_e(\Theta^x_{J\leftarrow j})^{-1}.
\label{eq:full-graph-dressed-edge}
\end{align}
For an action $\rho_x$ of $G$ on the channel-$x$ sample fiber, the root-framed child tuple
$\bar z^x_I=(\rho_x(\Theta^x_{I\leftarrow i})z^x_i)_{i\in I}$ collects the member coordinates in the root
frame, and a law $Q$ on the root sample fiber of $I$ is holonomy stabilized when
\begin{equation}
\widehat\rho_x(\mathsf k)Q=Q\qquad\text{for every }\mathsf k\in\mathrm{Hol}^x_I,
\label{eq:full-graph-holonomy-stabilization}
\end{equation}
which is law coherence up to holonomy stabilizers on the block (\Cref{def:geo-law-coherence}).

\propositionheading{Gauge laws of the rooted datum}{prop:rg-rooted-datum}
The map $p\mapsto\Theta^x_p$ on closed walks at $r_I$ is multiplicative under concatenation and invariant
under insertion or removal of backtracks, so it is a homomorphism on the fundamental group of
$\mathfrak I_I$ at $r_I$, with image $\mathrm{Hol}^x_I$ generated by the fundamental-cycle holonomies
$\mathrm{Hol}^x(p_e)$. Under a passive re-expression $g$ the tree transports, this homomorphism, and, for a
boundary copy $e:j\to i$ from $J$ to $I$, the link and the dressed edge transform as
\begin{align}
&\Theta^x_{I\leftarrow i}\longmapsto g_{r_I}\Theta^x_{I\leftarrow i}g_i^{-1},
\label{eq:rg-tree-transport-law}\\
&\pi_1(\mathfrak I_I,r_I)\longrightarrow G,\quad p\longmapsto\Theta^x_p,
\qquad
\Theta^x_p\longmapsto g_{r_I}\Theta^x_pg_{r_I}^{-1},
\label{eq:rg-full-holonomy-representation}\\
&\Theta^x_e\longmapsto g_i\Theta^x_eg_j^{-1},
\qquad
\widehat\Theta^x_e\longmapsto g_{r_I}\widehat\Theta^x_eg_{r_J}^{-1},
\label{eq:rg-dressed-boundary-edge}
\end{align}
and $\bar z^x_I\mapsto\rho_x(g_{r_I})\bar z^x_I$, one root element acting on every entry. The datum
consisting of $\bar z^x_I$, the homomorphism \eqref{eq:rg-full-holonomy-representation}, and the dressed
edges is therefore root framed and equivariant under one root re-expression per block. The homomorphism
does not depend on the trees, and a change of tree multiplies each tree transport on the left by a based
holonomy. The tree links can be set to the identity by a passive re-expression unique up to one element
per block, after which each non-tree internal link equals the conjugated fundamental-cycle holonomy it
closes. \status{ESTABLISHED}

\noindent\emph{Proof.}
Substitution cancels interior factors of path products and gives the law of $\bar z^x_I$, and a backtrack
contributes $\Theta^x_e\Theta^x_{\bar e}=e$. Every internal copy satisfies
$\Theta^x_e=(\Theta^x_{I\leftarrow i})^{-1}\mathrm{Hol}^x(p_e)\Theta^x_{I\leftarrow j}$, with $\mathrm{Hol}^x(p_e)=e$
on tree edges, so the product along a closed walk at $r_I$ telescopes to a product of fundamental-cycle
holonomies and their inverses; walks and holonomy do not involve the trees, and for a second tree
$\widetilde\Theta^x_{I\leftarrow i}(\Theta^x_{I\leftarrow i})^{-1}$ is the product along a closed walk at
$r_I$. Finally $g_i=g_{r_I}\Theta^x_{I\leftarrow i}$ makes every tree transport, hence every tree link, the
identity, is forced up to the choice of $g_{r_I}$, and turns a non-tree link into
$g_{r_I}\mathrm{Hol}^x(p_e)g_{r_I}^{-1}$. $\square$

The last statement is the maximal-tree gauge fixing of lattice gauge theory \citep{Creutz1977}, in
the spanning-tree form of \citet[Sec.~IV A, Eq.~(4.1)]{DelcampDittrich2017PhaseDiagram}. The
measure tier retains the raw root-framed datum, since quotienting each holonomy separately would lose its
orientation relative to the root features and boundary legs, and an orbit space would need a quotient
theorem not supplied here.

\propositionheading{Two full-law treatments of holonomy for a candidate parent}{prop:rg-parent-holonomy-alternatives}
In the setting of \Cref{thm:rg-pointwise-parent-datum} two alternatives are available, which need not
exclude each other at different retained-coordinate or quotient levels. (i) Groupoid-covariant parent:
supply bimeasurable actions of a typed groupoid on the fine random space, on $\mathsf Z_A$, on the
observation space, and on the structural data; let the fine generative and recognition laws be covariant
and the selected posterior versions compatible with the actions; and require
$C_A(g\cdot y,g\cdot D)=C_A(y,D)$ for every admitted arrow $g$ and measurable $D\subseteq\mathsf Z_A$,
together with covariance of the evaluator for the actions on $(m_A,\xi_A)$ and on $(b_A,o,h_A)$. Then the
parent generative, posterior, and recognition laws are covariant. Invariance at one slice $(o,X)$ holds for
the arrows that fix $(o,X)$ and the selected posterior versions; invariance of $q^b_A$ and $q^m_A$
separately is only a projection consequence and does not replace these hypotheses. (ii) Retention: the
component roots, the based holonomy homomorphisms, and the dressed boundary edges of
\Cref{prop:rg-rooted-datum} are placed in $\mathsf H_A$, and their joint law is pushed forward with no
invariance claim. Neither alternative selects $I$, $A$, a partition, or a membership kernel.
\status{ESTABLISHED}

\noindent\emph{Proof.}
(i) Substituting the invariance of $C_A$ and the covariant fine laws into
\eqref{eq:rg-pointwise-parent-generative} and \eqref{eq:rg-pointwise-parent-triple} gives covariance of the
three parent laws, and evaluator covariance carries it to the evaluator tier; at a fixed slice only arrows
fixing $(o,X)$ and the versions map the slice to itself. (ii) defines the retained coordinates, and the
pushforward is that of \Cref{thm:rg-pointwise-parent-datum}. $\square$

\section{Linear feature tier and tree dependence}

\propositionheading{Linear cross-scale maps and their tree dependence}{prop:rg-linear-cross-scale}
Let $\rho^{\mathbb W}_x:G\to\GL(\mathbb W_x)$ be a linear representation on a feature fiber, choose weights $w_{Ii}\geq0$ with
$\sum_{i\in I}w_{Ii}=1$, and set
\begin{equation}
(\mathsf C_xz)_I=\sum_{i\in I}w_{Ii}\rho^{\mathbb W}_x(\Theta^x_{I\leftarrow i})z_i,
\qquad
(\mathsf S_x\hat z)_i=\rho^{\mathbb W}_x(\Theta^x_{I\leftarrow i})^{-1}\hat z_I
\label{eq:rg-linear-cross-scale-maps}
\end{equation}
for a fine feature $z$ and a coarse feature $\hat z=(\hat z_I)_I$, with
$\rho^{\mathbb W}_{x,f}(g)=\bigoplus_i\rho^{\mathbb W}_x(g_i)$ and $\rho^{\mathbb W}_{x,c}(g)=\bigoplus_I\rho^{\mathbb W}_x(g_{r_I})$. (i)
$\mathsf C_x\mathsf S_x=I$, the averaged feature $(\mathsf C_xz)_I$ is a weighted mean of the entries of the
root-framed tuple, which it compresses, and
$\mathsf C'_x=\rho^{\mathbb W}_{x,c}(g)\mathsf C_x\rho^{\mathbb W}_{x,f}(g)^{-1}$,
$\mathsf S'_x=\rho^{\mathbb W}_{x,f}(g)\mathsf S_x\rho^{\mathbb W}_{x,c}(g)^{-1}$. (ii) At nested depths $0\to1\to2$, with
$\Theta^{x,12}_{A\leftarrow I}$ the tree transport of the coarse vertex $I$ to the root of its depth-two
block $A$, the conditions
$\Theta^{x,02}_{A\leftarrow i}=\Theta^{x,12}_{A\leftarrow I}\Theta^{x,01}_{I\leftarrow i}$ and
$w^{02}_{Ai}=w^{12}_{AI}w^{01}_{Ii}$ give $\mathsf C^{02}_x=\mathsf C^{12}_x\mathsf C^{01}_x$ and
$\mathsf S^{02}_x=\mathsf S^{01}_x\mathsf S^{12}_x$; nested forests alone give neither equality. (iii) For a
fine message operator $\mathcal A_x$ with $\mathcal A'_x=\rho^{\mathbb W}_{x,f}(g)\mathcal A_x\rho^{\mathbb W}_{x,f}(g)^{-1}$, the
fine-to-meta, meta-to-fine, and meta-to-meta operators
$\mathcal A^x_{I\leftarrow j}=(\mathsf C_x\mathcal A_x)_{Ij}$,
$\mathcal A^x_{j\leftarrow I}=(\mathcal A_x\mathsf S_x)_{jI}$, and
$\mathcal A^x_{I\leftarrow J}=(\mathsf C_x\mathcal A_x\mathsf S_x)_{IJ}$ transform between the corresponding
endpoint frames, and for a linear cross map $\Phi_f:\mathbb W_{b,f}\to \mathbb W_{m,f}$ with
$\Phi'_f=\rho^{\mathbb W}_{m,f}(g)\Phi_f\rho^{\mathbb W}_{b,f}(g)^{-1}$ the compressed map $\Phi^c=\mathsf C_m\Phi_f\mathsf S_b$
satisfies $\Phi^{c\prime}=\rho^{\mathbb W}_{m,c}(g)\Phi^c\rho^{\mathbb W}_{b,c}(g)^{-1}$. (iv) A second family of trees gives
transports $\widetilde\Theta^x_{I\leftarrow i}=\mathrm{Hol}^x(\gamma_i)\Theta^x_{I\leftarrow i}$, with
$\gamma_i$ the closed walk from $r_I$ to $i$ along the first tree and back along the second, and
\begin{equation}
\begin{gathered}
(\widetilde{\mathsf C}_xz)_I-(\mathsf C_xz)_I
=\sum_{i\in I}w_{Ii}\bigl(\rho^{\mathbb W}_x(\mathrm{Hol}^x(\gamma_i))-I\bigr)\rho^{\mathbb W}_x(\Theta^x_{I\leftarrow i})z_i,\\
(\widetilde{\mathsf S}_x\hat z)_i=\rho^{\mathbb W}_x(\Theta^x_{I\leftarrow i})^{-1}\rho^{\mathbb W}_x(\mathrm{Hol}^x(\gamma_i))^{-1}\hat z_I.
\end{gathered}
\label{eq:rg-linear-cross-scale-tree-change}
\end{equation}
The difference is not a root re-expression, so $\mathsf C_x$, $\mathsf S_x$, and the message operators
depend on the trees through gauge-invariant quantities: for two vertices joined by two edge copies with
$\mathrm{SO}(2)$ links the identity and the rotation by $\vartheta$, the defining representation, equal
weights, equal features $e_1$, and root at the first vertex, the two trees give averaged features of norms
$1$ and $\lvert\cos(\vartheta/2)\rvert$, although these norms are invariant under every root
re-expression. Tree independence holds for inputs whose root-framed features lie in the fixed space
$\operatorname{Fix}(\rho^{\mathbb W}_x\circ\mathrm{Hol}^x_I)=\{v\in \mathbb W_x:\rho^{\mathbb W}_x(\mathsf k)v=v\text{ for every }\mathsf k\in\mathrm{Hol}^x_I\}$,
for $\mathsf S_x$ on coarse features in that subspace, and for all inputs when
$\rho^{\mathbb W}_x\circ\mathrm{Hol}^x_I$ is trivial. \status{ESTABLISHED}

\noindent\emph{Proof.}
(i) $(\mathsf C_x\mathsf S_x\hat z)_I=\sum_iw_{Ii}\hat z_I=\hat z_I$, and substituting
\eqref{eq:rg-tree-transport-law} gives the covariance laws. (ii) Substitute the two conditions into the
composites. (iii) follows from (i). (iv) The closed walk $\gamma_i$ has path product
$\widetilde\Theta^x_{I\leftarrow i}(\Theta^x_{I\leftarrow i})^{-1}$, which gives
\eqref{eq:rg-linear-cross-scale-tree-change}. In the example the two trees give $e_1$ and
$\tfrac12(e_1+R_\vartheta e_1)$, of norm $\lvert\cos(\vartheta/2)\rvert$. If every
$\rho^{\mathbb W}_x(\Theta^x_{I\leftarrow i})z_i$ is fixed by the represented holonomy, the difference vanishes, and
likewise for $\mathsf S_x$. $\square$

This fixed space is the feature-level counterpart of \eqref{eq:full-graph-holonomy-stabilization}; a tree
choice is a gauge fixing for the retained datum but part of the model for the averaged feature
$\mathsf C_xz$.

\propositionheading{Composite trees and one rooted forest compose the tree transports}{prop:rg-forest-tree-composition}
At nested depths $0\to1\to2$ as in \Cref{prop:rg-linear-cross-scale}(ii), fix a channel $x$ with links
$\Theta^x_e\in G$, let the blocks $I$ of depth one partition $V$, let the depth-one multigraph $\mathfrak I_1$
of \Cref{def:full-graph-coarse-links} carry the dressed marks $\Theta^{x,1}_e=\widehat\Theta^x_e$ as links,
and let the blocks $A$ of depth two partition the vertices of $\mathfrak I_1$, every block connected; $A$ also
denotes the set of agents of its children $I\subseteq A$, with internal multigraph $\mathfrak I_A$ in
$\mathfrak I$. The transports $\Theta^{x,01}_{I\leftarrow i}$, $\Theta^{x,12}_{A\leftarrow I}$, and
$\Theta^{x,02}_{A\leftarrow i}$ are taken along rooted spanning trees of $\mathfrak I_I$, of the internal
multigraph of $A$ in $\mathfrak I_1$, and of $\mathfrak I_A$. (i) For any such trees of the children, with
roots $r_I$, and of $A$ in $\mathfrak I_1$, with root child $I_A$, the copies of these trees form a spanning
tree of $\mathfrak I_A$, the composite tree of $A$. With the composite tree and the root $r_{I_A}$, every agent
$i\in I\subseteq A$ and every copy $e:j\to i$ with $j\in J\subseteq A'$ satisfy
\begin{equation}
\begin{gathered}
\Theta^{x,02}_{A\leftarrow i}=\Theta^{x,12}_{A\leftarrow I}\Theta^{x,01}_{I\leftarrow i},\\
\Theta^{x,12}_{A\leftarrow I}\widehat\Theta^x_e\bigl(\Theta^{x,12}_{A'\leftarrow J}\bigr)^{-1}
=\Theta^{x,02}_{A\leftarrow i}\Theta^x_e\bigl(\Theta^{x,02}_{A'\leftarrow j}\bigr)^{-1},
\end{gathered}
\label{eq:rg-forest-tree-composition}
\end{equation}
so the depth-two link obtained by lifting twice with \eqref{eq:full-graph-coarse-links} is the one-step lift
through the composite trees: for $A'\neq A$ a dressed edge \eqref{eq:full-graph-dressed-edge}, and for $A'=A$
the fundamental-cycle holonomy $\mathrm{Hol}^x(p_e)$ of \eqref{eq:full-graph-tree-transport-law} for the
composite tree, which is $e_G$ on its copies. The transport condition of
\Cref{prop:rg-linear-cross-scale}(ii) therefore holds, and with the weights
$w^{02}_{Ai}=w^{12}_{AI}w^{01}_{Ii}$, which sum to one over $A$, both equalities of that part hold. (ii)
Choose, before the partitions and independently of them, a rooted spanning forest of $\mathfrak I$, that is,
a rooted spanning tree of each connected component of $\mathfrak I$, and let every block be a connected
subtree of it, meaning that the forest copies joining agents of the block connect the block. Each block then
has exactly one agent nearest to the root of its component, measured by the number of copies on the forest
path; call it $r_I$ for a child and $r_A$ for a block of depth two. Give each child the forest copies joining
its agents, rooted at $r_I$, and give $A$ in $\mathfrak I_1$ the forest copies joining distinct children of
$A$, rooted at the child $I_A$ that contains $r_A$. These are rooted spanning trees, $r_{I_A}=r_A$, and the
composite tree of $A$ consists of the forest copies joining agents of $A$, so (i) applies to the forest tree
of $A$ rooted at $r_A$, and
\begin{equation}
\Theta^{x,12}_{A\leftarrow I}=\Theta^{x,02}_{A\leftarrow r_I},
\qquad
\Theta^{x,02}_{A\leftarrow i}=\Theta^{x,02}_{A\leftarrow r_I}\Theta^{x,01}_{I\leftarrow i}.
\label{eq:rg-forest-root-factorization}
\end{equation}
For a longer chain of nested partitions whose blocks are all connected subtrees of the forest, the trees
induced at every depth $\ell$, with the agents as the blocks of depth zero, formed by the forest copies
inside a block of depth $\ell+1$ that join distinct blocks of depth $\ell$ and rooted at the block of depth
$\ell$ that contains its nearest agent, satisfy the hypotheses of (i) at the depths $\ell\to\ell+1\to\ell+2$
for every $\ell$. (iii) If the links of channel $x$ are graph flat on every block of depth two
(\Cref{def:geo-graph-flat-links}), in particular if they are frame aligned,
\eqref{eq:rg-forest-tree-composition} holds with any rooted spanning tree of $\mathfrak I_A$ with root
$r_{I_A}$ in place of the composite tree. The transports in (i) and (ii) are path products of the
group-valued links, whose reversed copies carry the inverse links; a matrix communication
map $\mathsf L_e$ of \Cref{def:geo-selective-matrix-link}, whose reversed copy is specified separately, is
not a transport, and on the vector tier, where $\rho_x$ acts linearly on the sample fibers, its dressings by
the represented transports of (i) satisfy
\begin{equation}
\begin{aligned}
&\rho_x(\Theta^{x,12}_{A\leftarrow I})\bigl[\rho_x(\Theta^{x,01}_{I\leftarrow i})\mathsf L_e
\rho_x(\Theta^{x,01}_{J\leftarrow j})^{-1}\bigr]\rho_x(\Theta^{x,12}_{A'\leftarrow J})^{-1}\\
&\qquad=\rho_x(\Theta^{x,02}_{A\leftarrow i})\mathsf L_e\rho_x(\Theta^{x,02}_{A'\leftarrow j})^{-1},
\end{aligned}
\label{eq:rg-forest-matrix-dressing}
\end{equation}
which inverts only represented transports, so $\mathsf L_e$ may be singular. \status{ESTABLISHED}

\noindent\emph{Proof.}
Deleting a backtrack leaves a path product unchanged, since $\Theta^x_e\Theta^x_{\bar e}=e_G$, and a walk in
a tree reduces in this way to the tree path between its endpoints, because a closed walk of positive length
in a tree backtracks at a vertex farthest from its base point. The argument for (i) uses only that the links
are elements of $G$ whose reversed copies carry the inverses, which holds at every depth by
\Cref{prop:full-graph-coarse-links}. (i) Counting reversal pairs, the trees of the children contain
$\lvert A\rvert$ less the number of children, each pair joining agents of one child, and the tree of $A$ in
$\mathfrak I_1$ contains the number of children less one, each pair joining distinct children; together
they connect the agents of $A$ with $\lvert A\rvert-1$ pairs, so they form a spanning tree. Let the tree path
$p$ from $I$ to $I_A$ in $\mathfrak I_1$ traverse the copies $e_a:j_a\to i_a$ from $I_{a-1}$ to $I_a$ for
$a=1,\ldots,\lvert p\rvert$, oriented along the path, with $I_0=I$, $I_{\lvert p\rvert}=I_A$, and $i_0=i$.
By \eqref{eq:full-graph-coarse-links},
\[
\Theta^{x,12}_{A\leftarrow I}\Theta^{x,01}_{I\leftarrow i}
=\Theta^{x,01}_{I_{\lvert p\rvert}\leftarrow i_{\lvert p\rvert}}\Theta^x_{e_{\lvert p\rvert}}
\bigl[(\Theta^{x,01}_{I_{\lvert p\rvert-1}\leftarrow j_{\lvert p\rvert}})^{-1}
\Theta^{x,01}_{I_{\lvert p\rvert-1}\leftarrow i_{\lvert p\rvert-1}}\bigr]
\Theta^x_{e_{\lvert p\rvert-1}}\cdots\Theta^x_{e_1}
\bigl[(\Theta^{x,01}_{I_0\leftarrow j_1})^{-1}\Theta^{x,01}_{I_0\leftarrow i_0}\bigr].
\]
Each bracket is the path product of a walk in the tree of $I_a$ from $i_a$ through $r_{I_a}$ to $j_{a+1}$, and
the leftmost factor is that of the tree path from $i_{\lvert p\rvert}$ to $r_{I_A}$, so the right side is the
path product of a walk in the composite tree from $i$ to $r_{I_A}$, which is $\Theta^{x,02}_{A\leftarrow i}$.
Substituting this at both endpoints of
$\widehat\Theta^x_e=\Theta^{x,01}_{I\leftarrow i}\Theta^x_e(\Theta^{x,01}_{J\leftarrow j})^{-1}$ gives the
second equality, whose right side is \eqref{eq:full-graph-dressed-edge} for $A'\neq A$ and
\eqref{eq:full-graph-tree-transport-law} for $A'=A$, equal to $e_G$ on tree copies as in
\Cref{prop:full-graph-coarse-links}. The weights are nonnegative with
$\sum_{I\subseteq A}w^{12}_{AI}\sum_{i\in I}w^{01}_{Ii}=1$, so \Cref{prop:rg-linear-cross-scale}(ii)
applies. (ii) If two distinct agents of a block were nearest to the root, the forest path between them would
lie in the block, since the forest copies inside the block connect it and paths in a tree are unique, and it
would pass through their last common ancestor, which is nearer to the root. The forest copies inside a block
connect it and contain no cycle, so they form a spanning tree of its internal multigraph. Inside $A$ lie
$\lvert A\rvert-1$ forest pairs, of which $\lvert A\rvert$ less the number of children lie inside children,
so the number of children less one join distinct children; these pairs connect the children,
because the forest copies inside $A$ connect $A$, so they form a spanning tree of $A$ in $\mathfrak I_1$, and
with the trees of the children they make up the forest copies inside $A$. The agent $r_A$ lies in
$I_A\subseteq A$ and is nearest in $A$, hence in $I_A$, so $r_{I_A}=r_A$. Then (i) applies; putting $i=r_I$,
where $\Theta^{x,01}_{I\leftarrow r_I}=e_G$, gives the first equality of
\eqref{eq:rg-forest-root-factorization}, and substituting it into (i) gives the second. For a longer chain the
same count shows that each induced tree is a spanning tree. For a block of depth $\ell+2$, its induced tree at
depth $\ell$ is the union of the induced trees at depth $\ell$ of its children and of the forest copies
joining distinct children, which form its induced tree at depth $\ell+1$, and the block of depth $\ell$ that
contains its nearest agent is the root at depth $\ell$ of the child containing that agent, whose nearest agent
it is. (iii) Graph-flat links have path product $e_G$ on every closed walk in the block
(\Cref{prop:geo-trivializing-criterion}), and another tree with the same root multiplies each transport on
the left by a based holonomy (\Cref{prop:rg-rooted-datum}), which leaves $\Theta^{x,02}_{A\leftarrow i}$
unchanged; frame-aligned links are graph flat (\Cref{rem:geo-flatness-strictness}). The matrix identity
follows from the first equality of \eqref{eq:rg-forest-tree-composition} at $i$ and at $j$ and the
multiplicativity of $\rho_x$. $\square$

Part (i) supplies what nesting alone leaves open in \Cref{prop:rg-linear-cross-scale}(ii): a depth-zero
tree of $A$ that contains the depth-one tree and is rooted at the root of the root child; another tree with
that root multiplies the direct transports $\Theta^{x,02}_{A\leftarrow i}$ by based holonomies
(\Cref{prop:rg-rooted-datum}). In (i) the trees may be chosen after each partition is realized, and the
composite tree then depends on the intermediate partition. In (ii) the forest is part of the blocking scheme
and is chosen before the partitions, so the direct transports of $A$ depend only on $A$ and the forest; a
random membership mark $R_\ell$ (\Cref{def:full-graph-membership-kernel}) enters (ii) only through realized
partitions whose blocks are connected subtrees of that forest, a realized block that is not such a subtree
is outside (ii), and overlapping membership kernels are outside the proposition, as in
\Cref{def:full-graph-coarse-links}.

Setting the tree links to the identity is the spanning-tree gauge fixing of lattice gauge theory in
the form of \citet[Sec.~IV A, Eq.~(4.1)]{DelcampDittrich2017PhaseDiagram}, whose change of spanning
tree at a fixed root is their Sec.~IV C, Eq.~(4.4); the coarse link as an ordered product of fine links
along a composite path is the macroscopic holonomy of
\citet[Sec.~4, Eq.~(4.1)]{Bahr2012OperatorSpinFoams}. Neither source composes trees across nested
partitions.

\section{Holonomy-conditioned marginal-law modes}
\label{sec:cg-holonomy-law-modes}

\definitionheading{Holonomy-conditioned marginal-law score}{def:cg-holonomy-law-score}
Fix a channel $x$, a connected block $I$ with root $r$, weights $w_i>0$ with $\sum_{i\in I}w_i=1$, and
marginal laws $q^x_i$ of the agents of $I$ on the sample fiber $\mathsf Z^x$, on which $G$ acts by
bimeasurable bijections $\rho_x(g)$ with pushforward $\widehat\rho_x(g)$. Let parent-law families
$\mathscr M^x_i\subseteq\operatorname{Prob}(\mathsf Z^x)$ be equivariant:
$\widehat\rho_x(\Theta^x_p)\mathscr M^x_j=\mathscr M^x_i$ for every walk $p$ of $\mathfrak I_I$ from $j$ to
$i$, and $\mathscr M^{x\prime}_i=\widehat\rho_x(g_i)\mathscr M^x_i$ under a passive re-expression. For walks
$p_i$ from $i$ to $r$ put $\bar q^x_i=\widehat\rho_x(\Theta^x_{p_i})q^x_i$, and with the based holonomy
group $\mathrm{Hol}^x_I(r)$ of \eqref{eq:full-graph-based-holonomy} define the fixed parent family and the
score
\begin{align}
\mathscr Q^x_{I,\mathrm{fix}}(r)&=\bigl\{Q\in\mathscr M^x_r:\widehat\rho_x(\mathsf k)Q=Q
\text{ for every }\mathsf k\in\mathrm{Hol}^x_I(r)\bigr\},
\label{eq:cg-holonomy-fixed-parent}\\
\mathfrak D^x_I&=\inf_{Q\in\mathscr Q^x_{I,\mathrm{fix}}(r)}\sum_{i\in I}w_i\KL(\bar q^x_i\Vert Q),
\qquad\inf\varnothing:=+\infty.
\label{eq:cg-holonomy-kl-score}
\end{align}
A minimizer, when one exists, is a holonomy-conditioned marginal-law mode.

\theoremheading{Invariant holonomy-conditioned marginal-law mode}{thm:cg-holonomy-kl-marginal}
The score $\mathfrak D^x_I$ does not depend on the walks, the root, or the passive frames. If the infimum
is attained, then $\mathfrak D^x_I=0$ exactly when some $Q_\star\in\mathscr Q^x_{I,\mathrm{fix}}(r)$
satisfies $\bar q^x_i=Q_\star$ for every $i\in I$, and the laws $\widehat\rho_x(\Theta^x_{p_i})^{-1}Q_\star$
then form a holonomy-stabilized parallel marginal-law section. If $\mathscr Q^x_{I,\mathrm{fix}}(r)=\varnothing$,
then $\mathfrak D^x_I=+\infty$. If $\mathfrak D^x_I=0$ without attainment, the conclusion is only forward-KL
approximability by fixed parent laws. \status{ESTABLISHED}

\noindent\emph{Proof.}
Relative entropy is invariant under bimeasurable bijections \eqref{eq:cg-kl-bimeasurable-invariance}.
Another walk $p'_i$ from $i$ to $r$ replaces $\bar q^x_i$ by $\widehat\rho_x(\mathsf k)\bar q^x_i$ with
$\mathsf k=\Theta^x_{p'_i}(\Theta^x_{p_i})^{-1}\in\mathrm{Hol}^x_I(r)$, and for every fixed parent $Q$,
$\KL(\widehat\rho_x(\mathsf k)\bar q^x_i\Vert Q)=\KL(\widehat\rho_x(\mathsf k)\bar q^x_i\Vert\widehat\rho_x(\mathsf k)Q)=\KL(\bar q^x_i\Vert Q)$,
so every objective value is path independent. A walk $\delta$ from $r$ to $r'$ conjugates
$\mathrm{Hol}^x_I(r)$ onto $\mathrm{Hol}^x_I(r')$, and pushforward by $\Theta^x_\delta$ maps the two fixed
families bijectively, by equivariance, and the transported laws of $r$ to those of $r'$, preserving every
term. A passive re-expression sends $\Theta^x_p$ to $g_i\Theta^x_pg_j^{-1}$ for $p$ from $j$ to $i$ and
$q^x_i$ to $\widehat\rho_x(g_i)q^x_i$, so it pushes every transported law by $\widehat\rho_x(g_r)$,
conjugates the holonomy group by $g_r$, and carries the fixed family bijectively, and invariance of
relative entropy gives gauge invariance. Relative entropy vanishes exactly between equal laws and every
$w_i$ is positive, which gives the zero case; since $Q_\star$ is holonomy fixed, the reconstructed laws do
not depend on the walks and are law coherent by \Cref{prop:geo-three-agreements}. The empty case is the
convention, and without attainment the infimum supplies $Q_n$ with score below $1/n$ and no minimizer.
$\square$

\propositionheading{Compact holonomy: closed form and attainment}{prop:cg-holonomy-compact-closed-form}
In the setting of \Cref{thm:cg-holonomy-kl-marginal}, let the represented holonomy group
$\mathfrak H_r=\rho_x(\mathrm{Hol}^x_I(r))$ be a compact metrizable group acting jointly measurably on
$\mathsf Z^x$, with Haar probability $d\mathsf k$, and put
\begin{align}
\bar P&=\sum_{i\in I}w_i\bar q^x_i,
\qquad
\bar M(B)=\int_{\mathfrak H_r}\bar P(\mathsf k^{-1}B)d\mathsf k.
\label{eq:cg-haar-symmetrized-mixture}
\intertext{Then $\bar M$ is $\mathfrak H_r$-invariant, and every $\mathfrak H_r$-invariant law $Q$ satisfies,
in $[0,+\infty]$,}
\sum_{i\in I}w_i\KL(\bar q^x_i\Vert Q)&=\sum_{i\in I}w_i\KL(\bar q^x_i\Vert\bar M)+\KL(\bar M\Vert Q).
\label{eq:cg-holonomy-compensation}
\end{align}
Consequently $\mathfrak D^x_I=\sum_iw_i\KL(\bar q^x_i\Vert\bar M)
+\inf_{Q\in\mathscr Q^x_{I,\mathrm{fix}}(r)}\KL(\bar M\Vert Q)$. If $\bar M\in\mathscr M^x_r$, for instance
when $\mathscr M^x_r=\operatorname{Prob}(\mathsf Z^x)$, the infimum is attained at $\bar M$, the score is the
holonomy-symmetrized information radius $\sum_iw_i\KL(\bar q^x_i\Vert\bar M)$, and $\bar M$ is the unique
minimizer when that value is finite. If instead $\mathsf Z^x=\R^K$, the holonomy acts linearly,
$\mathscr M^x_r$ is the nondegenerate Gaussian family, and $\bar M$ has finite second moments, mean
$\bar m$, nonsingular covariance $\bar S$, and $\KL(\bar M\Vert\mathcal N(\bar m,\bar S))<\infty$, then
$\mathcal N(\bar m,\bar S)$ is holonomy invariant and attains the infimum, uniquely when
$\mathfrak D^x_I<\infty$. For a finite holonomy group $d\mathsf k$ is uniform; a holonomy group that is only
relatively compact in a group acting jointly continuously on a Polish fiber has the same fixed laws as its
closure, to which the statement applies; and when every $\bar q^x_i$ is a nondegenerate Gaussian and the
holonomy acts linearly with compact closure, the Gaussian hypotheses hold and the score is finite.
\status{ESTABLISHED}

\noindent\emph{Proof.}
Left invariance of $d\mathsf k$ gives invariance of $\bar M$. Since $d\bar q^x_i/d\bar P\leq w_i^{-1}$,
splitting $\log(d\bar q^x_i/dQ)$ through $\bar P$ gives, for every $Q$, the compensation identity
$\sum_iw_i\KL(\bar q^x_i\Vert Q)=\sum_iw_i\KL(\bar q^x_i\Vert\bar P)+\KL(\bar P\Vert Q)$ in $[0,+\infty]$: the
negative part of the second term is integrable because $t(\log t)^-\leq e^{-1}$, and both sides are
infinite unless $\bar P\ll Q$. Applied with $Q$ and with $\bar M$, it reduces
\eqref{eq:cg-holonomy-compensation} to $\KL(\bar P\Vert Q)=\KL(\bar P\Vert\bar M)+\KL(\bar M\Vert Q)$ for
invariant $Q$, where both sides are infinite unless $\bar P\ll Q$. If $\bar P\ll Q$ with density $f$,
invariance of $Q$ gives $\bar M$ the invariant density $\bar f(z)=\int f(\mathsf k^{-1}z)d\mathsf k$; the
invariant set $\{\bar f=0\}$ is null for $\bar M$ and hence for $\bar P$, so $\bar P\ll\bar M\ll Q$. Averaging
the invariant function $\log\bar f$ over the group by Tonelli's theorem gives
$\int\log\bar fd\bar P=\int\log\bar fd\bar M=\KL(\bar M\Vert Q)$ with finite negative parts, and integrating
$\log f=\log(f/\bar f)+\log\bar f$ against $\bar P$ proves the identity. Attainment and uniqueness follow
because relative entropy vanishes only between equal laws. In the Gaussian case, invariance of $\bar M$ and
linearity give $\mathsf k\bar m=\bar m$ and $\mathsf k\bar S\mathsf k^\top=\bar S$, and
$\log[d\mathcal N(\bar m,\bar S)/d\mathcal N(\mu,\Sigma)]$ is a quadratic polynomial whose $\bar M$- and
$\mathcal N(\bar m,\bar S)$-expectations agree, which gives
$\KL(\bar M\Vert\mathcal N(\mu,\Sigma))=\KL(\bar M\Vert\mathcal N(\bar m,\bar S))+\KL(\mathcal N(\bar m,\bar S)\Vert\mathcal N(\mu,\Sigma))$.
For a relatively compact group, $\mathsf k\mapsto\int\phi(\mathsf kz)Q(dz)$ is continuous for bounded
continuous $\phi$ by dominated convergence, so fixed laws of the group and of its closure agree. For
Gaussian constituents with compact linear holonomy, $\bar M$ is a mixture of Gaussians with covariances in
a compact set of positive-definite matrices, so its density is bounded and its second moments are finite,
and convexity of relative entropy in its second argument bounds each $\KL(\bar q^x_i\Vert\bar M)$ by an
average of Gaussian relative entropies continuous on the compact group. $\square$

\corollaryheading{Noncompact Gaussian holonomy leaves no admissible parent}{cor:cg-holonomy-noncompact-gaussian}
Let $\mathsf Z^x=\R^K$, let the holonomy act by invertible linear maps, and let $\mathscr M^x_r$ be the
nondegenerate Gaussian family. Then $\mathscr Q^x_{I,\mathrm{fix}}(r)\neq\varnothing$ if and only if the
represented holonomy group is conjugate into $\mathrm O(K)$. If one loop holonomy has an eigenvalue off the
unit circle, for instance $\lvert\det\rvert\neq1$, or is not diagonalizable, then $\mathfrak D^x_I=+\infty$
for every configuration of marginals. For a block containing a loop with independent $\GL(K)$ links, the
link configurations with $\lvert\det\rvert=1$ on that loop form a Lebesgue-null set, so the score is
$+\infty$ off a null set of configurations; it distinguishes configurations only after the links are
restricted to a group conjugate into $\mathrm O(K)$ or a parent family with a nonempty fixed set is
declared. \status{ESTABLISHED}

\noindent\emph{Proof.}
$\mathcal N(\mu,\Sigma)$ with $\Sigma\succ0$ is fixed by $\mathsf k$ exactly when $\mathsf k\mu=\mu$ and
$\mathsf k\Sigma\mathsf k^\top=\Sigma$. If one such $\Sigma$ is fixed by the whole group, then
$\Sigma^{-1/2}\mathsf k\Sigma^{1/2}$ is orthogonal for every element; conversely, if
$S^{-1}\mathsf kS\in\mathrm O(K)$ for every element, then $\mathcal N(0,SS^\top)$ is fixed. An element
conjugate into $\mathrm O(K)$ is diagonalizable with eigenvalues of modulus one, and
$\det(\mathsf k\Sigma\mathsf k^\top)=(\det\mathsf k)^2\det\Sigma$; an empty fixed family gives $+\infty$ by
convention. The holonomy of a simple loop is a product of link matrices in which one chosen link $\Theta$
enters once, as $A_1\Theta^{\pm1}A_2$ with $A_1,A_2$ fixed by the other links; this is a diffeomorphism of
$\GL(K)$, so the preimage of the null set $\{\lvert\det\rvert=1\}$ is null, and Fubini's theorem over the
remaining links completes the argument. $\square$

With latent links (\Cref{ch:links}), the generators carry a law, and the score of
\Cref{def:cg-holonomy-law-score} becomes a function of them.

\corollaryheading{Latent noncompact links empty the Gaussian parent family almost surely}{cor:link-holonomy-parent-almost-sure}
Let \(G=\GL^+(K)\) with \(K\geq1\) act on the sample fiber \(\mathsf Z^x=\R^K\) of a channel \(x\) with latent
links by its defining representation, let the parent family be the nondegenerate Gaussian family as in
\Cref{cor:cg-holonomy-noncompact-gaussian}, and let \(I\) be a connected block, as in
\Cref{def:cg-holonomy-law-score}, whose internal multigraph
\(\mathfrak I_I\) is not a forest, that is, has a loop copy, two parallel copies, or a cycle through three or
more agents. Then \(\mathfrak I_I\) has a closed walk \(p=(e_0,\ldots,e_{r-1})\) that traverses some reversal
pair a different number of times in its two orientations, and for every such walk the following hold.
(i) For every choice of frames,
\(\det\mathrm{Hol}^x(p)=\exp\bigl(\sum_{a=0}^{r-1}\Tr\mathring\omega^x_{e_a}\bigr)\), with
\(\mathring\omega^x_{\bar e}=-\mathring\omega^x_e\).
(ii) The generator families with \(\sum_a\Tr\mathring\omega^x_{e_a}=0\) form a hyperplane of \(\mathscr T\), which
is Lebesgue-null, and for every generator family off this hyperplane the fixed parent family
\eqref{eq:cg-holonomy-fixed-parent} is empty and \(\mathfrak D^x_I=+\infty\) for every configuration of
marginals.
(iii) Under every law of the generators that is absolutely continuous with respect to Lebesgue measure on
\(\mathscr T\), in particular the generator prior of \Cref{def:link-generator-prior} whenever it has a
density, as it does when its base law is Gaussian with positive-definite covariances, and every recognition
law absolutely continuous with respect to such a prior, \(\mathfrak D^x_I=+\infty\) almost surely, and its
expectation is \(+\infty\).
(iv) A law of the generators under which \(\mathfrak D^x_I\) is finite with positive probability charges the
hyperplane of (ii); the point mass at \(\mathring\omega=0\) is such a law for nondegenerate Gaussian
marginals.
(v) For \(G=\mathrm{SL}(2,\R)\) acting on \(\R^2\) by its defining representation and a block whose internal
multigraph has cycle rank one, the generator families at which the represented holonomy group is conjugate
into \(\mathrm{SO}(2)\) contain a nonempty open subset of \(\mathscr T\); on it the fixed parent family is
nonempty, and every law of the generators with an everywhere positive density gives it positive
probability. \status{ESTABLISHED}

\noindent\emph{Proof.}
If \(\mathfrak I_I\) is a forest, every copy is a bridge, and a closed walk crosses a bridge as often in one
orientation as in the other, because it returns to the side it started from. A loop copy traversed once,
the walk \((e,\bar e')\) through two parallel copies \(e,e'\), and a cycle traversed once each traverse some
reversal pair once in one orientation and never in the other.
(i) By \eqref{eq:geo-edge-twist-holonomy},
\(\mathrm{Hol}^x(p)=U_{i_0}e^{\mathring\omega^x_{e_0}}\cdots e^{\mathring\omega^x_{e_{r-1}}}U_{i_0}^{-1}\), and
\(\det e^{\mathbf B}=e^{\Tr\mathbf B}\) by Jacobi's formula.
(ii) Collecting terms, \(\sum_a\Tr\mathring\omega^x_{e_a}=\sum_{e\in E_+}\mathrm n_e\Tr\mathring\omega^x_e\), where
\(\mathrm n_e\) is the number of traversals of \(e\) minus that of \(\bar e\). Some \(\mathrm n_e\) is nonzero,
and the family with \(\mathring\omega^x_e=I_K\) and every other generator zero gives the value
\(\mathrm n_eK\neq0\), so the functional is a nonzero linear form and its zero set is a hyperplane. Off the
hyperplane \(\lvert\det\mathrm{Hol}^x(p)\rvert\neq1\) by (i). For the root \(r\) of \(I\), the closed walk at
\(r\) that follows a walk from \(r\) to \(i_0\), then \(p\), then the reverse walk has a holonomy conjugate to
\(\mathrm{Hol}^x(p)\), with the same determinant, so a loop holonomy of the block has an eigenvalue off the
unit circle, and \Cref{cor:cg-holonomy-noncompact-gaussian} gives an empty fixed family and
\(\mathfrak D^x_I=+\infty\) for every configuration of marginals.
(iii) Such a law gives the null hyperplane probability zero, and a law absolutely continuous with respect
to one that is absolutely continuous with respect to Lebesgue measure is itself absolutely continuous with
respect to Lebesgue measure. The generator prior \(Z_{\mathscr S}^{-1}e^{-\mathscr S}\nu^{\mathscr T}\) of
\Cref{def:link-generator-prior} is absolutely continuous with respect to its base law, which has a Lebesgue
density when it is Gaussian with positive-definite covariances.
(iv) The first statement is (ii). At \(\mathring\omega=0\) every twist is the identity, so every holonomy is
the identity by \eqref{eq:geo-edge-twist-holonomy}; the trivial group is compact, and
\Cref{prop:cg-holonomy-compact-closed-form} gives a finite score for nondegenerate Gaussian marginals.
(v) With cycle rank one, \(\mathrm{Hol}^x_I\) is generated by one fundamental-cycle holonomy \(\mathsf k\)
(\Cref{prop:rg-rooted-datum}), which by \eqref{eq:geo-edge-twist-holonomy} is conjugate to an ordered product
of exponentials of the generators along the cycle, so \(\Tr\mathsf k\) is a continuous function of
\(\mathring\omega\). At the generator family with \(\frac\pi2\mathfrak J\) on the non-tree copy,
\(\mathfrak J=\begin{psmallmatrix}0&-1\\1&0\end{psmallmatrix}\), and zero elsewhere, \(\mathsf k\) is conjugate to
the rotation by \(\frac\pi2\), so the set where \(\lvert\Tr\mathsf k\rvert<2\) is open and nonempty. There
\(\mathsf k\in\mathrm{SL}(2,\R)\) has the eigenvalues \(e^{\pm\mathrm i\phi}\) with \(\sin\phi\neq0\); the real and
imaginary parts of an eigenvector for \(e^{\mathrm i\phi}\) form a basis of \(\R^2\) in which \(\mathsf k\) acts as
a rotation, so the group generated by \(\mathsf k\) is conjugate into \(\mathrm{SO}(2)\), and
\Cref{cor:cg-holonomy-noncompact-gaussian} gives a nonempty fixed family. A positive density has positive
integral over a nonempty open set. $\square$

The null-set sentence of \Cref{cor:cg-holonomy-noncompact-gaussian} measures link configurations in
\(\GL(K)\). With latent generators the exceptional set is the explicit hyperplane of (ii), linear in the
generators because \(\log\det e^{\mathbf B}=\Tr\mathbf B\), and it depends neither on the frames, nor on the
trace-free components of the generators, nor on whether they commute. The score is then not a statistic of
the recognition law: it is \(+\infty\) under every law with a density, and only a law that charges the
hyperplane, such as a point mass at a held link family, can make it finite. By (v), on
\(\mathrm{SL}(2,\R)\), which has no determinant direction, a set of generator families of positive
probability keeps a nonempty fixed family.

\remarkheading{Holonomy score on a three-agent loop}{rem:cg-holonomy-score-loop}
On the triangle of the running example (\Cref{sec:geo-running-example}), now with belief sample fiber
$\R^2$ and links in $\mathrm{SO}(2)$ acting by rotation, root the block at the first agent, take the tree
edges joining it to the other two, and let the loop holonomy be a rotation by $\phi$ with $\sin\phi\neq0$.
For transported marginals $\bar q^b_i=\mathcal N(m_i,\sigma^2I_2)$ with weights $w_i$ and the nondegenerate
Gaussian parent family, the fixed family is $\{\mathcal N(0,sI_2):s>0\}$ and
$\mathfrak D^b_I=\log(1+\tfrac{1}{2\sigma^2}\sum_iw_i\lvert m_i\rvert^2)$, attained at
$s=\sigma^2+\tfrac12\sum_iw_i\lvert m_i\rvert^2$, the moment-matched Gaussian of the Haar-symmetrized
mixture; rotations preserve norms, so the score does not depend on the walks. With frame-aligned links on
the same triangle the holonomy is trivial, every nondegenerate Gaussian is fixed, and the score is
$\tfrac12\log\det(I_2+\operatorname{Cov}_w(m)/\sigma^2)$, which never exceeds the value above. Replacing one
link by a matrix of determinant $2$ empties the fixed family and makes the score $+\infty$
(\Cref{cor:cg-holonomy-noncompact-gaussian}). \status{ESTABLISHED}

\noindent\emph{Proof.}
A fixed Gaussian has mean fixed by $R_\phi$, hence zero because $R_\phi-I$ is invertible, and a symmetric
covariance commuting with $R_\phi$; the linear conditions on its off-diagonal entry and the difference of
its diagonal entries have determinant $2\sin^2\phi\neq0$, so the covariance is scalar. In dimension two,
$\KL(\mathcal N(m,\sigma^2I_2)\Vert\mathcal N(0,sI_2))=\sigma^2/s+\lvert m\rvert^2/(2s)-1+\log(s/\sigma^2)$,
whose weighted sum is minimized at the stated $s$ with the stated value. In the flat case moment matching
gives $\mathcal N(\bar m,\sigma^2I_2+\operatorname{Cov}_w(m))$ and the stated value. With $c_1,c_2$ the
eigenvalues of $\operatorname{Cov}_w(m)/\sigma^2$ and
$\sum_iw_i\lvert m_i\rvert^2=\lvert\bar m\rvert^2+\operatorname{tr}\operatorname{Cov}_w(m)$, the inequality of
arithmetic and geometric means gives $\tfrac12\log((1+c_1)(1+c_2))\leq\log(1+\tfrac12(c_1+c_2))$, which is at
most the holonomy value. A link of determinant $2$ gives the loop holonomy a determinant of modulus $2$.
$\square$

\definitionheading{Exact common recovery}{def:cg-common-recovery}
A statistical experiment $\{P_\theta\}_{\theta\in\Theta}$ admits exact common recovery through a
parameter-independent Markov kernel $K_\downarrow$ when a parameter-independent Markov kernel $K_\uparrow$
satisfies
\begin{equation}
P_\theta K_\downarrow K_\uparrow=P_\theta\qquad\text{for every }\theta\in\Theta.
\label{eq:cg-holonomy-common-recovery}
\end{equation}

\remarkheading{A zero score is not common recovery}{rem:cg-marginal-not-joint}
The score \eqref{eq:cg-holonomy-kl-score} depends on the laws of the block only through their transported
marginals. When both marginals of a pair of agents are non-Dirac, there are joint laws $P_1\neq P_2$ with
the same marginals (\Cref{prop:prob-marginals-do-not-determine-joint}), and every kernel $K_\downarrow$ for
which $PK_\downarrow$ depends on $P$ only through its marginals sends them to one law, so the experiment
$\{P_1,P_2\}$ admits no common recovery through it. \status{ESTABLISHED}

\noindent\emph{Proof.}
If $P_1K_\downarrow=P_2K_\downarrow$, a kernel $K_\uparrow$ recovering both would give
$P_1=P_1K_\downarrow K_\uparrow=P_2K_\downarrow K_\uparrow=P_2$. $\square$

\corollaryheading{Parallel cross morphisms preserve fixed marginal-law sections}{cor:cg-holonomy-cross-morphism}
Let $\Phi:\mathcal E_b\to\mathcal E_m$ be a cross map, expressed in the frame of agent $i$ by
$\Phi_i:\mathcal B_b\to\mathcal B_m$, whose link-intertwining defects \eqref{eq:rg-cross-morphism-defects}
vanish on the internal copies of the block, so that
$\widehat\rho_m(\Theta^m_p)\circ\Phi_j=\Phi_i\circ\widehat\rho_b(\Theta^b_p)$ for every walk $p$ from $j$ to
$i$, and assume $\Phi_r(\mathscr M^b_r)\subseteq\mathscr M^m_r$. Then
$\Phi_r(\mathscr Q^b_{I,\mathrm{fix}}(r))\subseteq\mathscr Q^m_{I,\mathrm{fix}}(r)$, and applying $\Phi_i$
agentwise to a holonomy-stabilized parallel belief-law section gives such a model-law section. If
$\widetilde\Phi:\mathcal E_m\to\mathcal E_b$ likewise has vanishing link-intertwining defects and
$\widetilde\Phi_r(\mathscr M^m_r)\subseteq\mathscr M^b_r$, then
$\widetilde\Phi_r(\mathscr Q^m_{I,\mathrm{fix}}(r))\subseteq\mathscr Q^b_{I,\mathrm{fix}}(r)$, with the
analogous statement for sections. \status{ESTABLISHED}

\noindent\emph{Proof.}
Vanishing defects on the edges give the intertwining for walks by composition, using the coordinate form of
\Cref{prop:rg-defect-covariance}. If $Q\in\mathscr Q^b_{I,\mathrm{fix}}(r)$ and $p$ is a closed walk at
$r$, then $\widehat\rho_m(\Theta^m_p)\Phi_r(Q)=\Phi_r(\widehat\rho_b(\Theta^b_p)Q)=\Phi_r(Q)$, and family
compatibility gives the first inclusion. If the belief marginals transport to $Q$, the same square shows
that their images transport to $\Phi_r(Q)$. Interchanging the channels and replacing $\Phi$ by
$\widetilde\Phi$ gives the rest. $\square$

\section{Structural fixed sectors}
\label{sec:cg-partial}

\definitionheading{Connection Laplacian and fixed fiber}{def:cg-connection-laplacian}
Fix a channel and a block $I$ with internal multigraph $\mathfrak I_I$, let the channel's sample fiber be
a real vector space $\mathbb V$ of dimension $K$ with a linear representation $\rho$, and write $\Theta_e$ for the
represented link $\rho(\Theta^x_e)\in\GL(\mathbb V)$. For each unoriented internal copy choose an orientation
$e:j\to i$ and a positive-semidefinite form $W_e$ on $\mathbb V$. The connection Laplacian of the block is the
positive-semidefinite form $L_I$ on $\mathbb V^I$ with
\begin{equation}
z^\top L_Iz=\sum_{e:j\to i}(z_i-\Theta_ez_j)^\top W_e(z_i-\Theta_ez_j),
\label{eq:cg-connection-laplacian-energy}
\end{equation}
in which loops and parallel copies each contribute one term. Its fixed fiber and rank are
$\mathbb V_{I,\mathrm{fix}}=\ker L_I$ and $f_I=\dim\mathbb V_{I,\mathrm{fix}}$. For a root $r$ the represented based holonomy group
$\mathfrak H_r=\rho(\mathrm{Hol}^x_I(r))$ has the fixed space
$\operatorname{Fix}(\mathfrak H_r)=\{v\in \mathbb V:\mathsf kv=v\text{ for every }\mathsf k\in\mathfrak H_r\}$.

\propositionheading{Kernel-holonomy fixed-space isomorphism}{prop:cg-kernel-holonomy}
(i) The energy \eqref{eq:cg-connection-laplacian-energy} does not depend on the chosen orientations when
the reversed copy carries $W_{\bar e}=\Theta_e^\top W_e\Theta_e$. For positive-semidefinite weights,
$z\in\ker L_I$ exactly when
\begin{equation}
W_e^{1/2}(z_i-\Theta_ez_j)=0\qquad\text{for every internal copy }e:j\to i,
\label{eq:cg-psd-section-constraint}
\end{equation}
and $\ker L_I$ is the space $H^0$ of global sections of the cellular sheaf with vertex stalks $\mathbb V$ and
restriction maps $W_e^{1/2}$ at $i$ and $W_e^{1/2}\Theta_e$ at $j$, of which $L_I$ is the degree-zero
Laplacian \citep[Sec.~3.2 and Thm.~3.1]{HansenGhrist2019}; for orthogonal transports and scalar weights
$L_I$ is the vector-diffusion operator of \citet[Sec.~3]{SingerWu2012}. A passive re-expression, with the
weights re-expressed as $W_e\mapsto\rho(g_i)^{-\top}W_e\rho(g_i)^{-1}$, carries $\ker L_I$ to
$\rho_f(g)\ker L_I$, with $\rho_f(g)=\bigoplus_{i\in I}\rho(g_i)$. (ii) If $W_e\succ0$ on every internal
copy, then on every connected component $I_\alpha$ with root $r_\alpha$ root evaluation
$z\mapsto z_{r_\alpha}$ is a linear isomorphism
$\ker L_{I_\alpha}\to\operatorname{Fix}(\mathfrak H_{r_\alpha})$, so
$\ker L_I\cong\bigoplus_\alpha\operatorname{Fix}(\mathfrak H_{r_\alpha})$ and
$f_I=\sum_\alpha\dim\operatorname{Fix}(\mathfrak H_{r_\alpha})$. (iii) With positive-definite weights a
connected block has $f_I=K$ exactly when every represented holonomy element is the identity, equivalently
when the represented links are graph flat on $I$ (\Cref{prop:geo-trivializing-criterion}); this is the
full fixed sector, $0<f_I<K$ is a partial fixed sector, and $f_I=0$ leaves no retained degree. A
nonfaithful representation can make the represented holonomy trivial when the principal holonomy is not.
(iv) With positive-semidefinite weights holonomy may survive in weight-null directions: on a triangle with
$W_{12}=\operatorname{diag}(1,0)$, $\Theta_{12}=\operatorname{diag}(1,2)$, $W_{23}=\operatorname{diag}(0,1)$,
$\Theta_{23}=\operatorname{diag}(2,1)$, and $W_{31}=\Theta_{31}=I$, the loop product is $2I$ and fixes no
nonzero vector, while $\ker L_I$ is the two-dimensional space of constant configurations.
\status{ESTABLISHED}

\noindent\emph{Proof.}
(i) Substituting $\Theta_{\bar e}=\Theta_e^{-1}$ and $W_{\bar e}$ reproduces the term of $e$. Since
$L_I\succeq0$, $z\in\ker L_I$ exactly when $z^\top L_Iz=\sum_e\lvert W_e^{1/2}(z_i-\Theta_ez_j)\rvert^2=0$. The
diagonal blocks of $L_I$ are sums of $W_e$ and $\Theta_e^\top W_e\Theta_e$ and its off-diagonal blocks are
$-W_e\Theta_e$, the blocks of the sheaf Laplacian with the stated restriction maps, whose kernel is $H^0$ by
the discrete Hodge theorem. The energy is unchanged under $z_i\mapsto\rho(g_i)z_i$,
$\Theta_e\mapsto\rho(g_i)\Theta_e\rho(g_j)^{-1}$, and the weight law, so
$L'_I=\rho_f(g)^{-\top}L_I\rho_f(g)^{-1}$. (ii) With positive-definite weights, $z\in\ker L_{I_\alpha}$ exactly
when $z_i=\Theta_ez_j$ on every copy. Such a section is determined by its root value, and transporting that
value around a closed walk returns it, so root evaluation is injective with image in
$\operatorname{Fix}(\mathfrak H_{r_\alpha})$. Conversely, for $u\in\operatorname{Fix}(\mathfrak H_{r_\alpha})$
put $z_v=\rho(\Theta^x_{I_\alpha\leftarrow v})^{-1}u$. Tree constraints hold by construction, and for a
non-tree copy $e:j\to i$ the represented fundamental-cycle holonomy
$\rho(\mathrm{Hol}^x(p_e))=\rho(\Theta^x_{I_\alpha\leftarrow i})\Theta_e\rho(\Theta^x_{I_\alpha\leftarrow j})^{-1}$
of \eqref{eq:full-graph-tree-transport-law} fixes $u$, which gives $\Theta_ez_j=z_i$. The Laplacian is the
direct sum over components. (iii) $\operatorname{Fix}(\mathfrak H_r)=\mathbb V$ exactly when every element is the
identity, the fundamental-cycle holonomies generate the group (\Cref{prop:rg-rooted-datum}), and elements
of the kernel of $\rho$ are represented by the identity. (iv) The loop product is
$\operatorname{diag}(1,2)\operatorname{diag}(2,1)=2I$. The copy $31$ forces $z_3=z_1$, the copy $12$ equal
first coordinates of $z_1$ and $z_2$, and the copy $23$ equal second coordinates of $z_2$ and $z_3$, so the
kernel consists of the constant configurations. $\square$

\propositionheading{Fixed-section compression}{prop:cg-fixed-section-compression}
Let $\Lambda$ be a positive-semidefinite form on the fine configuration space, for instance the connection
Laplacian of the network plus positive-semidefinite self terms. For a partition $\mathcal P$ into blocks
with positive-semidefinite internal weights, let $\iota_I$ be a full-column matrix whose range is
$\ker L_I$, and $\iota_{\mathcal P}=\operatorname{blkdiag}_{I\in\mathcal P}\iota_I$. Then
\begin{equation}
L_I\iota_I=0,
\qquad
\Lambda_c=\iota_{\mathcal P}^\top\Lambda\iota_{\mathcal P}\succeq0,
\qquad
\Lambda_c\succ0\Longleftrightarrow\operatorname{range}(\iota_{\mathcal P})\cap\ker\Lambda=\{0\}.
\label{eq:cg-partial-properness}
\end{equation}
Internal edges contribute nothing to $\Lambda_c$, and cut edges become the rectangular endpoint data of
\Cref{thm:cg-rectangular-endpoint-closure}. With positive-definite weights, the tree construction in the
proof of \Cref{prop:cg-kernel-holonomy} extends a basis of root values to one such matrix.
\status{ESTABLISHED}

\noindent\emph{Proof.}
The first identity holds by construction and the second by congruence; $y^\top\Lambda_cy=0$ exactly when
$\iota_{\mathcal P}y\in\ker\Lambda$, and $\iota_{\mathcal P}$ is injective. An internal term vanishes on
$\operatorname{range}\iota_I=\ker L_I$, and a cut term restricted to the ranges is the energy of
\Cref{thm:cg-rectangular-endpoint-closure} with the blocks of $\iota$ as injections. $\square$

\propositionheading{Nested fixed sections compose exactly}{prop:cg-nested-sections-compose}
Let a larger block $A$ be a union of first-level blocks with positive-semidefinite internal weights, and
let $\iota_{A\leftarrow\mathcal P}=\operatorname{blkdiag}_{I\subset A}\iota_I$ and
$L^{(1)}_A=\iota_{A\leftarrow\mathcal P}^\top L_A\iota_{A\leftarrow\mathcal P}$. Then
$\ker L_A=\iota_{A\leftarrow\mathcal P}\ker L^{(1)}_A$, and for every injection
$\iota^{(2)}_A:\mathbb V_{A,\mathrm{fix}}\hookrightarrow\bigoplus_{I\subset A}\mathbb V_{I,\mathrm{fix}}$ with range $\ker L^{(1)}_A$ the fine
fixed-section injection is $\iota^{\mathrm{fine}}_A=\iota_{A\leftarrow\mathcal P}\iota^{(2)}_A$.
\status{ESTABLISHED}

\noindent\emph{Proof.}
A zero mode of $L_A$ has zero energy in every positive-semidefinite edge summand, so its restriction to each
first-level block lies in $\ker L_I$ and it equals $\iota_{A\leftarrow\mathcal P}y$ for some $y$. Since
$L_A\succeq0$, $y\in\ker L^{(1)}_A$ exactly when $y^\top L^{(1)}_Ay=0$, that is, exactly when
$\iota_{A\leftarrow\mathcal P}y\in\ker L_A$, which proves the first identity; substituting the range of
$\iota^{(2)}_A$ gives the second, and associativity of matrix composition iterates the statement.
$\square$

\remarkheading{Structural fixed sector versus marginal law occupation}{rem:cg-structural-vs-marginal}
Relative entropy between the marginals of two agents needs one of them transported into the frame of the
other and depends on the walk: a second walk changes the transported law by a loop-holonomy action, and
the two values agree for every pair of laws only under stabilization; with
$\mathsf k=\operatorname{diag}(1,-1,-1)\in\mathrm{SO}(3)$ and $q_i=q_j=\mathcal N(e_2,I_3)$, the identity walk
gives zero and a walk differing by $\mathsf k$ gives $2$. An unpenalized sum of nonnegative within-block
pairwise divergences vanishes at, and is minimized by, the singleton partition; it contains only marginal
information and does not test common recovery (\Cref{rem:cg-marginal-not-joint}). The structural fixed
sector and the marginal-law score are independent: on a two-vertex tree with identity transport and
$W=I_K$ the represented holonomy is trivial, yet $\KL(\mathcal N(-ae_1,I_K)\Vert\mathcal N(ae_1,I_K))=2a^2$ is
arbitrarily large; on a triangle with $K=3$, unit weights, identity transports on two edges, and closing
transport $\mathsf k$, the represented holonomy is nontrivial and
$\operatorname{Fix}(\mathsf k)=\operatorname{span}(e_1)$, so $f_I=1<K$, while $\mathsf k$ preserves
$\mathcal N(0,\sigma^2I_3)$ and assigning that law to every vertex makes every transported marginal
divergence zero. Transported divergence thus measures the occupation of the available law modes, not the
dimension of the structural fixed sector. \status{ESTABLISHED}

\noindent\emph{Proof.}
For equal covariances $\KL(\mathcal N(m_1,I)\Vert\mathcal N(m_2,I))=\tfrac12\lvert m_1-m_2\rvert^2$, which
gives $\tfrac12\lvert e_2-\mathsf ke_2\rvert^2=2$ and $2a^2$. A sum of nonnegative terms over pairs inside
blocks is zero for singleton blocks. The fixed space of $\mathsf k$ is $\operatorname{span}(e_1)$, and
$\mathsf k\sigma^2I_3\mathsf k^\top=\sigma^2I_3$. $\square$

\section{Cut links and rectangular sheaf closure}
\label{sec:cg-sheaf-closure}

Let $I$ and $J$ be blocks reduced to root coordinates $z_I,z_J\in \mathbb V$, with represented dressed boundary
edges $\widehat\Theta_e\in\GL(\mathbb V)$ (\Cref{def:rg-rooted-tree-data}) and positive-definite weights $W_e$ on
the cut copies $e$ from $J$ to $I$. The cut energy
$\mathcal E_{\mathrm{cut}}(z_I,z_J)=\sum_e(z_I-\widehat\Theta_ez_J)^\top W_e(z_I-\widehat\Theta_ez_J)$
equals $z_I^\top\Lambda_{II}z_I-2z_I^\top\Lambda_{IJ}z_J+z_J^\top\Lambda_{JJ}z_J$ with
$\Lambda_{II}=\sum_eW_e$, $\Lambda_{IJ}=\sum_eW_e\widehat\Theta_e$, and
$\Lambda_{JJ}=\sum_e\widehat\Theta_e^\top W_e\widehat\Theta_e$. Coarse families admit newly generated
positive-semidefinite self terms. The cut energy carries no factor one half, so $W_e$ and the blocks
$\Lambda_{II}$, $\Lambda_{IJ}$, $\Lambda_{JJ}$ are the coefficients of the quadratic form
$\mathcal E_{\mathrm{cut}}$, and a Gaussian law proportional to $\exp(-\mathcal E_{\mathrm{cut}})$ has
precision blocks $2\Lambda_{II}$, $-2\Lambda_{IJ}$, and $2\Lambda_{JJ}$. Supplement G writes every
quadratic energy with the factor one half, so its weights and precision blocks are twice the $W_e$
and $\Lambda$ blocks here. The excess, the admissible interval, and the balanced weight below are
homogeneous of degree one in the weights, and the links they determine of degree zero, so no
conclusion of this section depends on the convention.

\definitionheading{Anchored and no-new-anchor presentations of a cut}{def:cg-no-new-anchor-pair}
An anchored one-link presentation of the cut is a tuple $(\Theta,W,\Lambda^{\mathrm s}_I,\Lambda^{\mathrm s}_J)$
with $W\succ0$, $\Theta\in\GL(\mathbb V)$, and $\Lambda^{\mathrm s}_I,\Lambda^{\mathrm s}_J\succeq0$ such that, for all
$z_I,z_J$,
\begin{equation}
\mathcal E_{\mathrm{cut}}(z_I,z_J)=(z_I-\Theta z_J)^\top W(z_I-\Theta z_J)
+z_I^\top\Lambda^{\mathrm s}_Iz_I+z_J^\top\Lambda^{\mathrm s}_Jz_J.
\label{eq:cg-cut-anchored-link}
\end{equation}
The no-new-anchor pair subfamily consists of the presentations with
$\Lambda^{\mathrm s}_I=\Lambda^{\mathrm s}_J=0$: one link and one weight, with no new self terms.

\lemmaheading{Completed square for a cut}{lem:cg-cut-completed-square}
One has $\Lambda_{JJ}\succ0$ and, with $\bar\Theta=\Lambda_{II}^{-1}\Lambda_{IJ}$,
\begin{align}
\mathcal E_{\mathrm{cut}}(z_I,z_J)&=(z_I-\bar\Theta z_J)^\top\Lambda_{II}(z_I-\bar\Theta z_J)+z_J^\top\Delta z_J,
\label{eq:cg-cut-completed-square}\\
\Delta&=\Lambda_{JJ}-\Lambda_{IJ}^\top\Lambda_{II}^{-1}\Lambda_{IJ}
=\sum_e(\widehat\Theta_e-\bar\Theta)^\top W_e(\widehat\Theta_e-\bar\Theta)\succeq0.
\label{eq:cg-cut-excess}
\end{align}
(i) An anchored presentation exists if and only if $\det\Lambda_{IJ}\neq0$; with links restricted to
$\GL^+(\mathbb V)$ the condition is $\det\Lambda_{IJ}>0$. (ii) When $\det\Lambda_{IJ}\neq0$, the anchored presentations are exactly
$\Theta=W^{-1}\Lambda_{IJ}$, $\Lambda^{\mathrm s}_I=\Lambda_{II}-W$, and
$\Lambda^{\mathrm s}_J=\Lambda_{JJ}-\Lambda_{IJ}^\top W^{-1}\Lambda_{IJ}$, with $W$ in the Loewner interval
\begin{equation}
\mathfrak W_{\mathrm{cut}}=\{W:\Lambda_{IJ}\Lambda_{JJ}^{-1}\Lambda_{IJ}^\top\preceq W\preceq\Lambda_{II}\}
=\{W\succ0:W\preceq\Lambda_{II},\ \Lambda_{IJ}^\top W^{-1}\Lambda_{IJ}\preceq\Lambda_{JJ}\}.
\label{eq:cg-cut-admissible-weights}
\end{equation}
The choice $W=\Lambda_{II}$ gives $\Lambda^{\mathrm s}_I=0$ and $\Lambda^{\mathrm s}_J=\Delta$, and
$W=\Lambda_{IJ}\Lambda_{JJ}^{-1}\Lambda_{IJ}^\top$ gives $\Lambda^{\mathrm s}_J=0$, the reversed completion with
link $\Lambda_{JJ}^{-1}\Lambda_{IJ}^\top$ from $I$ to $J$ and weight $\Lambda_{JJ}$. (iii) When
$\det\Lambda_{IJ}\neq0$, the anchored presentation is unique if and only if $\Delta=0$, and it is then the
bare pair $(\bar\Theta,\Lambda_{II},0,0)$; otherwise there are infinitely many. The condition $\Delta=0$
itself forces $\det\Lambda_{IJ}\neq0$. (iv) The no-new-anchor pair subfamily contains a presentation of the
cut if and only if $\Delta=0$, that is, if and only if all cut twists $\widehat\Theta_e$ coincide. (v) The
weighted mean $\bar\Theta$ has the gauge law of a dressed edge,
$\bar\Theta\mapsto\rho(g_{r_I})\bar\Theta\rho(g_{r_J})^{-1}$, but it can be singular: equal weights with
twists $I$ and $-I$ give $\bar\Theta=0$. \status{ESTABLISHED}

\noindent\emph{Proof.}
Each $\widehat\Theta_e$ is invertible and $W_e\succ0$, so $\Lambda_{JJ}\succ0$. Completing the square in
$z_I$ gives \eqref{eq:cg-cut-completed-square}, and expanding the sum in \eqref{eq:cg-cut-excess} gives
$\Lambda_{JJ}-\Lambda_{IJ}^\top\Lambda_{II}^{-1}\Lambda_{IJ}$. Matching the three blocks of
\eqref{eq:cg-cut-anchored-link} gives $W\Theta=\Lambda_{IJ}$, $W+\Lambda^{\mathrm s}_I=\Lambda_{II}$, and
$\Theta^\top W\Theta+\Lambda^{\mathrm s}_J=\Lambda_{JJ}$, hence the formulas of (ii); both self terms are
positive semidefinite exactly on the second description in \eqref{eq:cg-cut-admissible-weights}. The link
is invertible exactly when $\Lambda_{IJ}$ is, and $\det\Theta=\det\Lambda_{IJ}/\det W$ has its sign. If
$\det\Lambda_{IJ}\neq0$, then $\Lambda_{IJ}^\top W^{-1}\Lambda_{IJ}\preceq\Lambda_{JJ}$ is equivalent to
$W^{-1}\preceq\Lambda_{IJ}^{-\top}\Lambda_{JJ}\Lambda_{IJ}^{-1}$ and, since inversion reverses the Loewner
order on positive-definite matrices, to $W\succeq\Lambda_{IJ}\Lambda_{JJ}^{-1}\Lambda_{IJ}^\top$; this is the
first description, which contains $W=\Lambda_{II}$ because $\Delta\succeq0$, and proves (i). The reversed
completion is the same computation with $z_I$ and $z_J$ exchanged. Expanding the right side shows
$\Lambda_{II}-\Lambda_{IJ}\Lambda_{JJ}^{-1}\Lambda_{IJ}^\top=\Lambda_{II}\Lambda_{IJ}^{-\top}\Delta\Lambda_{JJ}^{-1}\Lambda_{IJ}^\top$,
so the interval is a single point exactly when $\Delta=0$; otherwise every point of the segment between
its endpoints is admissible and distinct weights give distinct presentations. If $\Delta=0$ then
$\Lambda_{IJ}^\top\Lambda_{II}^{-1}\Lambda_{IJ}=\Lambda_{JJ}\succ0$, so $\Lambda_{IJ}$ is invertible. For (iv),
a presentation without self terms needs $W=\Lambda_{II}$, $W\Theta=\Lambda_{IJ}$, and
$\Theta^\top W\Theta=\Lambda_{JJ}$, that is, $\Delta=0$, and the sum in \eqref{eq:cg-cut-excess} with
$W_e\succ0$ vanishes exactly when every $\widehat\Theta_e=\bar\Theta$. For (v), invariance of the energy
under $z_I\mapsto\rho(g_{r_I})z_I$ and $z_J\mapsto\rho(g_{r_J})z_J$ transforms $W_e$ at the receiving end
and $\widehat\Theta_e$ as in \eqref{eq:rg-dressed-boundary-edge}, and substitution gives the law of
$\bar\Theta$; the example is direct. $\square$

\propositionheading{Balanced presentation of a cut}{prop:cg-balanced-cut-presentation}
In the setting of \Cref{lem:cg-cut-completed-square} with $\det\Lambda_{IJ}\neq0$, write
$W_{\mathrm{lo}}=\Lambda_{IJ}\Lambda_{JJ}^{-1}\Lambda_{IJ}^\top$ and $W_{\mathrm{hi}}=\Lambda_{II}$ for the endpoints
of $\mathfrak W_{\mathrm{cut}}$, and call an anchored presentation balanced when the self term at $J$,
carried to $I$ along the presentation's own link, equals the self term at $I$,
\begin{equation}
\Theta^{-\top}\Lambda^{\mathrm s}_J\Theta^{-1}=\Lambda^{\mathrm s}_I,
\qquad\text{equivalently}\qquad
\Theta^\top\Lambda_{II}\Theta=\Lambda_{JJ},
\label{eq:cg-cut-balance}
\end{equation}
so that the balanced link carries the form $\Lambda_{JJ}$ to the form $\Lambda_{II}$. (i) Exactly one
anchored presentation is balanced. Its weight is the matrix geometric mean of the endpoints,
\begin{equation}
W_{\mathrm{bal}}=W_{\mathrm{lo}}^{1/2}\bigl(W_{\mathrm{lo}}^{-1/2}W_{\mathrm{hi}}W_{\mathrm{lo}}^{-1/2}\bigr)^{1/2}W_{\mathrm{lo}}^{1/2},
\label{eq:cg-cut-balanced-weight}
\end{equation}
the unique positive-definite solution of $WW_{\mathrm{lo}}^{-1}W=W_{\mathrm{hi}}$, and its link is
$\Theta_{\mathrm{bal}}=W_{\mathrm{bal}}^{-1}\Lambda_{IJ}$. (ii) Under root re-expressions the weight transforms
as a weight at $I$, $W_{\mathrm{bal}}\mapsto\rho(g_{r_I})^{-\top}W_{\mathrm{bal}}\rho(g_{r_I})^{-1}$, the link as a
dressed edge, $\Theta_{\mathrm{bal}}\mapsto\rho(g_{r_I})\Theta_{\mathrm{bal}}\rho(g_{r_J})^{-1}$, and each self
term as a form at its own block. The presentation is end symmetric: read from $J$, with link
$\Theta_{\mathrm{bal}}^{-1}$ and weight $\Theta_{\mathrm{bal}}^\top W_{\mathrm{bal}}\Theta_{\mathrm{bal}}$, it is the
balanced presentation of the reversed cut, so that weight is the geometric mean of the endpoints
$\Lambda_{IJ}^\top\Lambda_{II}^{-1}\Lambda_{IJ}$ and $\Lambda_{JJ}$ of $J$. (iii) For scalar weights
$W_e=w_eI$ and orthogonal twists, an anchored presentation is balanced exactly when its link is
orthogonal; then $W_{\mathrm{bal}}=(\Lambda_{IJ}\Lambda_{IJ}^\top)^{1/2}$ and $\Theta_{\mathrm{bal}}$ is the
orthogonal polar factor of $\bar\Theta$. \status{ESTABLISHED}

\noindent\emph{Proof.}
By part (ii) of the lemma an anchored presentation has $\Theta=W^{-1}\Lambda_{IJ}$,
$\Lambda^{\mathrm s}_I=\Lambda_{II}-W$, and $\Lambda^{\mathrm s}_J=\Lambda_{JJ}-\Lambda_{IJ}^\top W^{-1}\Lambda_{IJ}$,
so $\Theta^{-\top}\Lambda^{\mathrm s}_J\Theta^{-1}=W\Lambda_{IJ}^{-\top}\Lambda_{JJ}\Lambda_{IJ}^{-1}W-W
=WW_{\mathrm{lo}}^{-1}W-W$, and the first form of \eqref{eq:cg-cut-balance} reads
$WW_{\mathrm{lo}}^{-1}W=W_{\mathrm{hi}}$. The second form is
$\Lambda^{\mathrm s}_J=\Theta^\top\Lambda^{\mathrm s}_I\Theta$ with the common term $\Theta^\top W\Theta$
removed from both sides, and after the substitution $\Theta=W^{-1}\Lambda_{IJ}$ it reads
$W^{-1}W_{\mathrm{hi}}W^{-1}=W_{\mathrm{lo}}^{-1}$, the same equation. (i) The substitution
$W=W_{\mathrm{lo}}^{1/2}YW_{\mathrm{lo}}^{1/2}$ is a bijection of positive-definite matrices and turns the
equation into $Y^2=M$ with $M=W_{\mathrm{lo}}^{-1/2}W_{\mathrm{hi}}W_{\mathrm{lo}}^{-1/2}\succ0$. A
positive-definite $Y$ with $Y^2=M$ commutes with $M$, so it is diagonal in an eigenbasis of $M$ with
positive eigenvalues whose squares are those of $M$; hence $Y=M^{1/2}$, which is
\eqref{eq:cg-cut-balanced-weight}. Since $\Delta\succeq0$ gives $W_{\mathrm{lo}}\preceq W_{\mathrm{hi}}$, one has
$M\succeq I$, hence $I\preceq M^{1/2}\preceq M$, and congruence by $W_{\mathrm{lo}}^{1/2}$ gives
$W_{\mathrm{lo}}\preceq W_{\mathrm{bal}}\preceq W_{\mathrm{hi}}$; the solution lies in $\mathfrak W_{\mathrm{cut}}$ and
is the weight of exactly one anchored presentation. (ii) As in the proof of part (v) of the lemma, a
root re-expression acts by $W_e\mapsto g_I^{-\top}W_eg_I^{-1}$ and
$\widehat\Theta_e\mapsto g_I\widehat\Theta_eg_J^{-1}$ with $g_I=\rho(g_{r_I})$ and $g_J=\rho(g_{r_J})$, so
$\Lambda_{II}$, $\Lambda_{IJ}$, and $\Lambda_{JJ}$ transform by $g_I^{-\top}(\cdot)g_I^{-1}$,
$g_I^{-\top}(\cdot)g_J^{-1}$, and $g_J^{-\top}(\cdot)g_J^{-1}$, and both endpoints by
$g_I^{-\top}(\cdot)g_I^{-1}$. If $W$ solves $WW_{\mathrm{lo}}^{-1}W=W_{\mathrm{hi}}$, then $g_I^{-\top}Wg_I^{-1}$
solves the transformed equation, so uniqueness gives the law of $W_{\mathrm{bal}}$, and substitution gives
the laws of the link and of the self terms. For end symmetry,
$(z_I-\Theta z_J)^\top W(z_I-\Theta z_J)=(z_J-\Theta^{-1}z_I)^\top\Theta^\top W\Theta(z_J-\Theta^{-1}z_I)$, so an
anchored presentation of the cut, read from $J$, is the anchored presentation
$(\Theta^{-1},\Theta^\top W\Theta,\Lambda^{\mathrm s}_J,\Lambda^{\mathrm s}_I)$ of the reversed cut, whose blocks
are $\Lambda_{JJ}$, $\Lambda_{IJ}^\top$, and $\Lambda_{II}$. The second form of \eqref{eq:cg-cut-balance} is
equivalent to $\Theta^{-\top}\Lambda_{JJ}\Theta^{-1}=\Lambda_{II}$, the balance condition of the reversed cut,
so the balanced presentation read from $J$ is balanced for the reversed cut, and (i) applied to that cut
identifies its weight. (iii) Here $\Lambda_{II}=\Lambda_{JJ}=(\sum_ew_e)I$, so the second form of
\eqref{eq:cg-cut-balance} says that $\Theta$ is orthogonal. For orthogonal $\Theta$ the weight
$W=\Lambda_{IJ}\Theta^\top$ is positive definite exactly when $\Lambda_{IJ}=W\Theta$ is a polar decomposition
of $\Lambda_{IJ}$, and then $W^2=\Lambda_{IJ}\Lambda_{IJ}^\top$ forces $W=(\Lambda_{IJ}\Lambda_{IJ}^\top)^{1/2}$ and
$\Theta=W^{-1}\Lambda_{IJ}$, the polar factor of $\Lambda_{IJ}$ and of $\bar\Theta=\Lambda_{IJ}/\sum_ew_e$.
$\square$

\propositionheading{No balanced presentation of a cut at the level of laws}{prop:cg-no-balanced-law-presentation}
Let $\operatorname{Gau}(\mathbb V)$ be the fiber of nondegenerate Gaussian laws on $\mathbb V$, and let $k$ copies $e$ from $J$ to
$I$ carry dressed edges $\widehat\Theta_e\in\GL(\mathbb V)$ and weights $w_e>0$ with $\sum_ew_e\widehat\Theta_e^{-1}$
invertible. For $q_J=\mathcal N(m_J,\Sigma_J)\in\operatorname{Gau}(\mathbb V)$ the transported laws
$(\widehat\Theta_e)_\#q_J=\mathcal N(\widehat\Theta_em_J,\Sigma_e)$, $\Sigma_e=\widehat\Theta_e\Sigma_J\widehat\Theta_e^\top$,
form the cut of \Cref{lem:cg-cut-completed-square} with weights $W_e=\tfrac12w_e\Sigma_e^{-1}$ and twists
$\widehat\Theta_e$; its blocks $\Lambda_{II}(\Sigma_J)$, $\Lambda_{IJ}(\Sigma_J)$, $\Lambda_{JJ}(\Sigma_J)$, its excess
$\Delta(\Sigma_J)$ \eqref{eq:cg-cut-excess}, and its balanced weight and link $W_{\mathrm{bal}}(\Sigma_J)$ and
$\Theta_{\mathrm{bal}}(\Sigma_J)$ \eqref{eq:cg-cut-balanced-weight} are functions of $\Sigma_J$, with
$\Lambda_{IJ}(\Sigma_J)$ invertible for every $\Sigma_J$. The cut has excess, $\Delta(\Sigma_J)\neq0$ for one and
hence for every $\Sigma_J$, exactly when the dressed edges are not all equal. Let $\mathscr C\subseteq\operatorname{Gau}(\mathbb V)$
be nonempty and closed under rescaling of the covariance, so that $\mathcal N(m,\Sigma)\in\mathscr C$ implies
$\mathcal N(m,t\Sigma)\in\mathscr C$ for every $t>0$. A presentation of the copies at the level of laws on
$\mathscr C$ is a weight $W'>0$, an image map $\mathfrak T:\mathscr C\to\operatorname{Gau}(\mathbb V)$, and self terms $\mathfrak S_I:\operatorname{Gau}(\mathbb V)\to\R$
and $\mathfrak S_J:\mathscr C\to\R$ with
\begin{equation}
\sum_ew_e\KL\bigl(q_I\Vert(\widehat\Theta_e)_\#q_J\bigr)=W'\KL\bigl(q_I\Vert \mathfrak T[q_J]\bigr)+\mathfrak S_I(q_I)+\mathfrak S_J(q_J)
\quad\text{for all }q_I\in\operatorname{Gau}(\mathbb V)\text{ and }q_J\in\mathscr C,
\label{eq:cg-cut-law-presentation}
\end{equation}
and it is balanced when its Gaussian mean part is the link term of the balanced presentation of the cut of
$\Sigma_J$, that is, writing $\mathfrak T[q_J]=\mathcal N(\mu_{\mathfrak T},\Sigma_{\mathfrak T})$,
\begin{equation}
\tfrac12W'\Sigma_{\mathfrak T}^{-1}=W_{\mathrm{bal}}(\Sigma_J)\quad\text{and}\quad\mu_{\mathfrak T}=\Theta_{\mathrm{bal}}(\Sigma_J)m_J
\qquad\text{for every }q_J=\mathcal N(m_J,\Sigma_J)\in\mathscr C.
\label{eq:cg-cut-law-balance}
\end{equation}
If the cut has excess, no presentation at the level of laws on $\mathscr C$ is balanced. Two copies with dressed
edges $I$ and $2I$ and weights $w_1,w_2$ have $\Delta(\Sigma_J)=\frac{w_1w_2}{2(4w_1+w_2)}\Sigma_J^{-1}$, so they
admit no balanced presentation on any such $\mathscr C$. \status{ESTABLISHED}

\noindent\emph{Proof.}
Since $\Sigma_e^{-1}=\widehat\Theta_e^{-\top}\Sigma_J^{-1}\widehat\Theta_e^{-1}$, the blocks are
$\Lambda_{II}(\Sigma_J)=\tfrac12\sum_ew_e\widehat\Theta_e^{-\top}\Sigma_J^{-1}\widehat\Theta_e^{-1}$,
$\Lambda_{IJ}(\Sigma_J)=\tfrac12\bigl(\sum_ew_e\widehat\Theta_e^{-\top}\bigr)\Sigma_J^{-1}$, and
$\Lambda_{JJ}(\Sigma_J)=\tfrac12\bigl(\sum_ew_e\bigr)\Sigma_J^{-1}$, so $\Lambda_{IJ}(\Sigma_J)$ is invertible
exactly when $\sum_ew_e\widehat\Theta_e^{-1}$ is, and each block is homogeneous of degree $-1$ in $\Sigma_J$.
Hence so are the endpoints $W_{\mathrm{lo}}=\Lambda_{IJ}\Lambda_{JJ}^{-1}\Lambda_{IJ}^\top$ and
$W_{\mathrm{hi}}=\Lambda_{II}$, the excess $\Delta$, and, since \eqref{eq:cg-cut-balanced-weight} is jointly
homogeneous of degree one in its two arguments, the balanced weight,
$W_{\mathrm{bal}}(t\Sigma_J)=t^{-1}W_{\mathrm{bal}}(\Sigma_J)$. By part (iv) of the lemma, $\Delta(\Sigma_J)=0$
exactly when all $\widehat\Theta_e$ coincide, a condition that does not involve $\Sigma_J$.

Fix a presentation, a law $q_J=\mathcal N(m_J,\Sigma_J)\in\mathscr C$, and $q_I=\mathcal N(m_I,\Sigma_I)\in\operatorname{Gau}(\mathbb V)$;
write $\mathfrak T[q_J]=\mathcal N(\mu_{\mathfrak T},\Sigma_{\mathfrak T})$, $B=\tfrac12W'\Sigma_{\mathfrak T}^{-1}$, and abbreviate
$\Lambda_{II}=\Lambda_{II}(\Sigma_J)$, $\Lambda_{IJ}=\Lambda_{IJ}(\Sigma_J)$. For Gaussian laws
$\KL(\mathcal N(m_1,\Sigma_1)\Vert\mathcal N(m_2,\Sigma_2))=\tfrac12(m_1-m_2)^\top\Sigma_2^{-1}(m_1-m_2)
+\tfrac12[\Tr(\Sigma_2^{-1}\Sigma_1)-\dim\mathbb V+\log\det\Sigma_2-\log\det\Sigma_1]$, so both sides of
\eqref{eq:cg-cut-law-presentation} are quadratic polynomials in $m_I$, and their difference is
\begin{align*}
0={}&m_I^\top(\Lambda_{II}-B)m_I-2m_I^\top(\Lambda_{IJ}m_J-B\mu_{\mathfrak T})+\Tr[(\Lambda_{II}-B)\Sigma_I]\\
&+\tfrac12\bigl(W'-\textstyle\sum_ew_e\bigr)\log\det\Sigma_I-\mathfrak S_I(q_I)+\mathrm c(q_J),
\end{align*}
where $\mathrm c(q_J)$ collects the terms that read $q_J$ alone. Subtracting this identity for a second law
$q_J'\in\mathscr C$ removes $\mathfrak S_I(q_I)$ and the $\log\det\Sigma_I$ term and leaves a quadratic polynomial in $m_I$
that vanishes for every $m_I\in\mathbb V$ and every $\Sigma_I\succ0$, so its quadratic and linear coefficients
vanish: $M=\Lambda_{II}(\Sigma_J)-B$ and $v=\Lambda_{IJ}(\Sigma_J)m_J-B\mu_{\mathfrak T}$ take the same value at every
$q_J\in\mathscr C$. If the presentation is balanced, \eqref{eq:cg-cut-law-balance} gives
$B=W_{\mathrm{bal}}(\Sigma_J)$ and
$B\mu_{\mathfrak T}=W_{\mathrm{bal}}(\Sigma_J)\Theta_{\mathrm{bal}}(\Sigma_J)m_J=\Lambda_{IJ}(\Sigma_J)m_J$, so $v=0$ and
$M=\Lambda_{II}(\Sigma_J)-W_{\mathrm{bal}}(\Sigma_J)$ for every $q_J\in\mathscr C$. Comparing $q_J$ with its
rescaling $\mathcal N(m_J,t\Sigma_J)\in\mathscr C$, the homogeneity gives $M=t^{-1}M$ for every $t>0$, hence
$M=0$ and $W_{\mathrm{bal}}(\Sigma_J)=W_{\mathrm{hi}}(\Sigma_J)$. Substituting into $WW_{\mathrm{lo}}^{-1}W=W_{\mathrm{hi}}$
gives $W_{\mathrm{hi}}W_{\mathrm{lo}}^{-1}W_{\mathrm{hi}}=W_{\mathrm{hi}}$, so $W_{\mathrm{lo}}=W_{\mathrm{hi}}$, and the
identity $\Lambda_{II}-\Lambda_{IJ}\Lambda_{JJ}^{-1}\Lambda_{IJ}^\top=\Lambda_{II}\Lambda_{IJ}^{-\top}\Delta\Lambda_{JJ}^{-1}\Lambda_{IJ}^\top$
of the proof of the lemma, whose outer factors are invertible, gives $\Delta(\Sigma_J)=0$: the cut has no
excess. For the two copies $I$ and $2I$, $\Sigma_1=\Sigma_J$ and $\Sigma_2=4\Sigma_J$, so
$\Lambda_{II}=\tfrac18(4w_1+w_2)\Sigma_J^{-1}$, $\Lambda_{IJ}=\tfrac14(2w_1+w_2)\Sigma_J^{-1}$, and
$\Lambda_{JJ}=\tfrac12(w_1+w_2)\Sigma_J^{-1}$, whence
$\Delta=\bigl[\tfrac12(w_1+w_2)-\tfrac{(2w_1+w_2)^2}{2(4w_1+w_2)}\bigr]\Sigma_J^{-1}=\frac{w_1w_2}{2(4w_1+w_2)}\Sigma_J^{-1}$.
$\square$

The balance condition is imposed only on the mean part of the link term, so the conclusion does not depend on
whether the image map is gauge covariant or the self terms are end symmetric: a self term that reads one law
cannot carry $\Lambda_{II}-W_{\mathrm{bal}}$, which is homogeneous of degree $-1$ in the covariance of the other
law, and the excess of the cut is exactly what keeps that difference from vanishing.

The exact closed category is larger than one invertible link; its endpoint maps and stalks use the notation
of the comparison maps and spaces of \Cref{ch:selective-communication}, which are an instance.

\theoremheading{Rectangular endpoint-map closure}{thm:cg-rectangular-endpoint-closure}
Choose a minimal factor $W_e=(S^Q_e)^\top S^Q_e$ with $S^Q_e:\mathbb V\to H_e$ of full row rank, and put
$S^K_e=S^Q_e\Theta_e$. Then
$(z_i-\Theta_ez_j)^\top W_e(z_i-\Theta_ez_j)=\lVert S^Q_ez_i-S^K_ez_j\rVert^2_{H_e}$. For linear maps
$\iota_{iI}:\mathbb V_I\to \mathbb V$ from block coordinates, the coarse endpoint maps $S^{Q,c}_{I,e}=S^Q_e\iota_{iI}$ and
$S^{K,c}_{J,e}=S^K_e\iota_{jJ}$ reproduce the Galerkin compression of the cut energy exactly, parallel cut
copies compress losslessly into the direct sum of their stalks $H_e$, and nested blocking is associative.
Parallel copies between the same two coarse vertices also admit a minimal exact compression: with
coboundary rows $\delta_e=(S^{Q,c}_{I,e},-S^{K,c}_{J,e})$, one edge with stalk $H$ and coboundary $\delta$
reproduces their summed energy exactly when $\delta^\top\delta=\sum_e\delta_e^\top\delta_e$; the minimal stalk
dimension is $\operatorname{rank}\sum_e\delta_e^\top\delta_e\leq\dim \mathbb V_I+\dim \mathbb V_J$, attained by any
full-row-rank factor, and two full-row-rank factors differ by an orthogonal change of edge basis.
\status{ESTABLISHED}

\noindent\emph{Proof.}
Expanding the square gives the pair form. Precomposition gives the coarse maps, whose cross block
$-(S^{Q,c}_{I,e})^\top S^{K,c}_{J,e}=-\iota_{iI}^\top W_e\Theta_e\iota_{jJ}$ and diagonal blocks agree with the
compressed energy. Stacking in $\bigoplus_eH_e$ sums the energies, and a second injection $\iota'$ gives
$(S\iota)\iota'=S(\iota\iota')$. Two quadratic forms agree exactly when their symmetric matrices agree, which
gives the compression criterion and the minimal stalk dimension; for full-row-rank factors
$\delta_1,\delta_2$ of one matrix, $O=\delta_2\delta_1^+$ satisfies $O^\top O=I$ and $\delta_2=O\delta_1$.
$\square$

This is the degree-zero cellular-sheaf Laplacian category \citep{HansenGhrist2019}, in which vertex spaces
may have unequal dimensions and automorphism groups $\GL(\mathbb V_I)$ and edge basis changes act on both endpoint
maps and the edge metric. Twice the quadratic disagreement $V_e$ of \Cref{def:obs-selective-comparison} is
the edge term of such an energy once its comparison precision is absorbed into the endpoint maps, as in
the discourse sheaves of \citet{HansenGhrist2021Discourse}; the closure selects no decoder and no law, and
the normalized realization for fixed maps is \eqref{eq:obs-selective-attention-record}.

\propositionheading{No continuous extension of balanced cut links}{open:cg-cut-compression}
There is no rule assigning a matrix-valued link and a positive-semidefinite weight and self terms to
every cut that is continuous in the constituent weights and twists and agrees with
\Cref{prop:cg-balanced-cut-presentation} whenever $\Lambda_{IJ}$ is invertible. This already fails for two
parallel copies with identity weights and twists in $\mathrm{SO}(2)\subset\GL^+(2)$. Consequently adding
gauge covariance, invariance under reversal, or a rank requirement on singular links cannot produce such
a rule on all singular and rectangular cuts. This obstruction concerns extension of the balanced
presentation; the exact rectangular endpoint-map closure of \Cref{thm:cg-rectangular-endpoint-closure}
remains available. \status{ESTABLISHED}

\noindent\emph{Proof.}
For $-1<t<1$, take $W_1=W_2=I_2$ and
\begin{equation}
\widehat\Theta_1(t)=\begin{pmatrix}t&-\sqrt{1-t^2}\\\sqrt{1-t^2}&t\end{pmatrix},
\qquad \widehat\Theta_2(t)=\widehat\Theta_1(t)^\top.
\label{eq:cg-cut-continuous-extension-witness}
\end{equation}
The twists lie in $\mathrm{SO}(2)$ and depend continuously on $t$, while
$\Lambda_{II}=\Lambda_{JJ}=2I_2$ and $\Lambda_{IJ}=2tI_2$. For $t\neq0$, the interval endpoints are
$W_{\mathrm{lo}}=2t^2I_2$ and $W_{\mathrm{hi}}=2I_2$, so the balanced weight and link are
$W_{\mathrm{bal}}=2|t|I_2$ and $\Theta_{\mathrm{bal}}=\operatorname{sign}(t)I_2$.
The same limiting cut at $t=0$ is approached from both sides, but the links tend to $I_2$ and $-I_2$.
A matrix-valued continuous extension would have to equal both limits. $\square$

\section{The frame sector}
\label{sec:cg-frame}

\propositionheading{The frame sector}{prop:cg-frame-sector}
Fix a block $I$ whose constituent frames are $u_i=\sigma\cdot U_i^{-1}$ relative to a reference point
$\sigma$ (\Cref{prop:geo-frame-comparisons}), and a coarse datum $(u_I,q_I)$ with $q_I\in\mathcal B_x$ and
$u_I=\sigma\cdot U_I^{-1}$ a point of the same torsor, so that $u_I=u_i\cdot T_{iI}$ with
$T_{iI}=U_iU_I^{-1}$. (i) The constituent expressions
$\widehat\rho_x(T_{iI})q_I=\widehat\rho_x(U_i)\bar q_I$ of the coarse law, with
$\bar q_I=\widehat\rho_x(U_I)^{-1}q_I$ the coordinate in $\sigma$ of the point $[u_I,q_I]$ of the associated
fiber, depend on $(u_I,q_I)$ only through that point; for the Gaussian law fiber of the defining
representation of $\GL^+(K)$ the constituent expression of $q_I=\mathcal N(\mu_I,\Sigma_I)$ is
$\mathcal N(U_i\bar m,U_i\bar SU_i^\top)$ with $\bar m=U_I^{-1}\mu_I$ and $\bar S=U_I^{-1}\Sigma_IU_I^{-\top}$.
The data with the same constituent expressions form exactly one orbit
$\{(u_I\cdot g,\widehat\rho_x(g)^{-1}q_I):g\in G\}$, on which $G$ acts freely, so the coarse frame occupies a
copy of $G$, of dimension $K^2$ for $G=\GL^+(K)$, that the constituents do not see. If only the represented
frame $\rho_x(U_I)$ is recorded, a nonfaithful representation determines the principal coarse frame only
modulo its kernel. (ii) If $G$ contains an element $v\neq e$ with $v^n=e$, no map $F:G^n\to G$ is both left
equivariant, $F(gg_1,\ldots,gg_n)=gF(g_1,\ldots,g_n)$, and invariant under permutations of its arguments. For
$G=\GL^+(K)$ with $K\geq2$ and every $n\geq2$ such an element exists, a rotation of order $n$ in one
coordinate plane. The torsion hypothesis cannot be dropped: on the additive group $\R^d$ the arithmetic
mean is left equivariant and symmetric. \status{ESTABLISHED}

\noindent\emph{Proof.}
(i) By the quotient convention \eqref{eq:geo-quotient-convention},
$[u_I,q_I]=[u_i\cdot T_{iI},q_I]=[u_i,\widehat\rho_x(T_{iI})q_I]$, so the constituent expression is the
coordinate of $[u_I,q_I]$ in the frame $u_i$ and determines that point. Likewise
$[u_I,q_I]=[\sigma,\widehat\rho_x(U_I)^{-1}q_I]=[u_i\cdot U_i,\bar q_I]=[u_i,\widehat\rho_x(U_i)\bar q_I]$, and
the Gaussian pushforward by the linear map $U_i$ gives the Gaussian form. Two data define the same
point exactly when they differ by $(u_I,q_I)\mapsto(u_I\cdot g,\widehat\rho_x(g)^{-1}q_I)$, and
$u_I\cdot g=u_I$ forces $g=e$. Elements of the kernel of $\rho_x$ give equal represented frames. (ii) Left
multiplication by $v$ permutes the tuple $(e,v,\ldots,v^{n-1})$ cyclically, because $v\cdot v^{n-1}=e$.
Symmetry gives $F$ the same value at the permuted tuple, while equivariance multiplies it by $v$; hence
$vF=F$ and $v=e$, a contradiction. A rotation by $2\pi/n$ in one coordinate plane lies in $\GL^+(K)$ and has
order $n$. The mean on $\R^d$ satisfies $\tfrac1n\sum_i(g+g_i)=g+\tfrac1n\sum_ig_i$. $\square$

A rule $u_I=\sigma\cdot F(a_1,\ldots,a_n)$ with $u_i=\sigma\cdot a_i$ is independent of $\sigma$ exactly when
$F$ is left equivariant, because $\sigma'=\sigma\cdot g_0$ replaces every $a_i$ by $g_0^{-1}a_i$; (ii)
therefore excludes every symmetric rule of this kind for a meta-agent's frame.

\section{An operational record at \texorpdfstring{$c_*$}{c*} sensitive to loop holonomy}

\theoremheading{An operational record sensitive to network-loop holonomy}{thm:cg-loop-operational-holonomy}
Let $p$ be a closed walk of the interaction multigraph based at a probe agent $i_0$ whose belief and model
sample fibers are real vector spaces $\mathbb V_b,\mathbb V_m$ with linear representations $\rho_b,\rho_m$, and write
$z_{i_0}(y)\in \mathbb V_{\mathrm{pr}}=\mathbb V_b\oplus \mathbb V_m$ for the probe coordinates of a latent configuration $y$ and
$\rho_{\mathrm{pr}}=\rho_b\oplus\rho_m$. Let $\Omega$ be a finite set of labels $\omega$ with prior masses
$\psi_\omega>0$, each prescribing link values on the edges of $p$, with represented loop holonomy
$\mathrm{Hol}^{\mathrm{pr}}_\omega=\rho_b(\mathrm{Hol}^b_\omega(p))\oplus\rho_m(\mathrm{Hol}^m_\omega(p))$. Keep
the joint $\mathbb P(dy,do)$ independent of $\omega$, with selected posterior version
$o\mapsto\boldsymbol\Pi_o$, the slice $\boldsymbol\Pi_{\theta,V,o,X}$ at fixed parameters and structural
data; fix $\Sigma_{\mathrm r}\succ0$ on $\mathbb V_{\mathrm{pr}}$, and declare the record joint
\[
P^{\mathrm{rec}}(\omega,dy,do,dr)=\psi_\omega\mathbb P(dy,do)
\mathcal N\bigl(r;(\mathrm{Hol}^{\mathrm{pr}}_\omega-I)z_{i_0}(y),\Sigma_{\mathrm r}\bigr)dr.
\]
(i) $P^{\mathrm{rec}}$ is a normalized joint. At an admitted $o$ and given $\omega$, the record is
distributed as $(\mathrm{Hol}^{\mathrm{pr}}_\omega-I)Z+\epsilon$, with $Z\sim(z_{i_0})_\#\boldsymbol\Pi_o$ and
$\epsilon\sim\mathcal N(0,\Sigma_{\mathrm r})$ independent of $Z$. Two labels $\omega,\omega'$ have the same
record law exactly when $(\mathrm{Hol}^{\mathrm{pr}}_\omega-I)Z$ and $(\mathrm{Hol}^{\mathrm{pr}}_{\omega'}-I)Z$
have the same law. (ii) If $\mathrm{Hol}^{\mathrm{pr}}_{\omega_0}=I$ and $\mathrm{Hol}^{\mathrm{pr}}_{\omega_1}-I$
is invertible, as for $\mathrm{SO}(2)$ links in both channels whose loop rotation angles are not multiples
of $2\pi$, the record laws of $\omega_0$ and $\omega_1$ differ exactly when $(z_{i_0})_\#\boldsymbol\Pi_o$ is
not the point mass at $0$. Under the same condition the statistic $\mathsf q=r^\top\Sigma_{\mathrm r}^{-1}r$
has different laws: its mean is $\dim \mathbb V_{\mathrm{pr}}$ under $\omega_0$ and
$\dim \mathbb V_{\mathrm{pr}}+\E\bigl[Z^\top(\mathrm{Hol}^{\mathrm{pr}}_{\omega_1}-I)^\top\Sigma_{\mathrm r}^{-1}
(\mathrm{Hol}^{\mathrm{pr}}_{\omega_1}-I)Z\bigr]\in(\dim \mathbb V_{\mathrm{pr}},+\infty]$ under $\omega_1$. (iii)
Under a passive re-expression $g$, with the record expressed in the probe frame,
$r\mapsto\rho_{\mathrm{pr}}(g_{i_0})r$ and
$\Sigma_{\mathrm r}\mapsto\rho_{\mathrm{pr}}(g_{i_0})\Sigma_{\mathrm r}\rho_{\mathrm{pr}}(g_{i_0})^\top$, the
record joint is the image of the original under a bimeasurable bijection; the posterior over $\omega$ given
$(o,r)$ and the law of $\mathsf q$ are unchanged. \status{ESTABLISHED}

\noindent\emph{Proof.}
(i) The Gaussian factor integrates to one in $r$, and $\psi$ and $\mathbb P$ are normalized. The label is
independent of $(y,o)$, so given $\omega$ and $o$ the configuration has law $\boldsymbol\Pi_o$ and the
record kernel is $\mathcal N((\mathrm{Hol}^{\mathrm{pr}}_\omega-I)z_{i_0}(y),\Sigma_{\mathrm r})$. The
characteristic function of the record is
$\E[e^{is^\top(\mathrm{Hol}^{\mathrm{pr}}_\omega-I)Z}]e^{-s^\top\Sigma_{\mathrm r}s/2}$, whose Gaussian factor
never vanishes, so two record laws agree exactly when the laws of the two images of $Z$ agree. (ii) Under
$\omega_0$ the image is $0$. Under $\omega_1$ invertibility gives
$(\mathrm{Hol}^{\mathrm{pr}}_{\omega_1}-I)Z=0$ almost surely exactly when $Z=0$ almost surely; a rotation by
$\alpha$ in the plane has $\det(R_\alpha-I)=2-2\cos\alpha$. Given $Z$, the conditional mean of $\mathsf q$ is
$\lvert\Sigma_{\mathrm r}^{-1/2}(\mathrm{Hol}^{\mathrm{pr}}_\omega-I)Z\rvert^2+\dim \mathbb V_{\mathrm{pr}}$, and
Tonelli's theorem gives the two means; the second exceeds the first exactly when the image of $Z$ is
nonzero with positive probability. (iii) The loop product telescopes under
$\Theta^x_e\mapsto g_i\Theta^x_eg_j^{-1}$, so with $U=\rho_{\mathrm{pr}}(g_{i_0})$ one has
$\mathrm{Hol}^{\mathrm{pr}}_\omega\mapsto U\mathrm{Hol}^{\mathrm{pr}}_\omega U^{-1}$ and
$z_{i_0}\mapsto Uz_{i_0}$, hence the new mean is $U(\mathrm{Hol}^{\mathrm{pr}}_\omega-I)z_{i_0}$ and the new
record kernel is the image of the old one under $r\mapsto Ur$. The joint is thus the image under the
bijection that changes the latent coordinates and multiplies $r$ by $U$, which leaves the posterior over
$\omega$ unchanged, and $r^\top U^\top(U\Sigma_{\mathrm r}U^\top)^{-1}Ur=r^\top\Sigma_{\mathrm r}^{-1}r$.
$\square$

\paragraph{Physicist's summary.}
A block of agents at one base point is a small lattice gauge system whose invariants are loop products and
edge twists modulo one global conjugation. A coarse description keeps a root frame, the tree transports to
it, the holonomy of every independent loop, and the tree-dressed boundary links; the tree is a gauge
fixing, as in lattice gauge theory, but averaged member features depend on it whenever the holonomy acts
on them. The configurations that survive exact aggregation form the kernel of the connection Laplacian,
for positive-definite couplings the holonomy-fixed vectors, and the members' marginal laws have one parent
exactly when their transported laws coincide in a holonomy-invariant law, so a planar rotation loop forces
a zero parent mean and a volume-stretching loop allows no Gaussian parent. Links between blocks compress
exactly into rectangular endpoint maps, and into one link with new self-energies when an explicit matrix is
invertible, uniquely once the self-energies are required to balance across the link; at the level of the
laws themselves no such balanced form exists once the copies disagree. The parent frame is
pure gauge for the constituents, and a noisy record of
how the loop rotates the probe's state distinguishes holonomies unless that state is pinned at zero;
Supplement B computes that holonomy from a connection along a base loop.
